\section{บทที่ 4 อนุกรมและทฤษฎีอนุกรมฟูเรียร์}

\textbf{ชื่อบทภาษาอังกฤษ:} Series and Fourier Series Theory\\
\textbf{รายวิชา:} คณิตศาสตร์วิศวกรรมไฟฟ้า (Electrical Engineering Mathematics)\\
\textbf{รหัสวิชา:} 5571102\\
\textbf{สาขาวิชา:} เทคโนโลยีวิศวกรรมไฟฟ้า\\
\textbf{คณะ:} เทคโนโลยีอุตสาหกรรม\\
\textbf{ระดับผู้เรียน:} นักศึกษาระดับปริญญาตรี\\
\textbf{ช่วงเวลาตามแผนการสอน:} สัปดาห์ที่ 6--7\\
\textbf{จำนวนชั่วโมงที่แนะนำ:} บรรยายและกิจกรรมคำนวณรวม 6 ชั่วโมง\\
\textbf{ความยาวเป้าหมาย:} 25--30 หน้า A4\\
\textbf{รูปแบบการเรียน:} บรรยาย การคำนวณแบบมีส่วนร่วม การอภิปราย และการวิเคราะห์กราฟ\\
\textbf{โหมดการออกแบบ:} Topic-Based Design Mode

%\par\noindent\rule{\textwidth}{0.4pt}\par

\subsection{1. แนวคิดสำคัญของบท}

สัญญาณไฟฟ้าที่พบในงานจริงมิได้มีรูปเป็นไซน์สมบูรณ์เสมอไป แรงดันจากวงจรสวิตชิ่ง สัญญาณพัลส์จากระบบดิจิทัล กระแสของโหลดไม่เชิงเส้น และสัญญาณควบคุมแบบมอดูเลตความกว้างพัลส์ล้วนเป็นรูปคลื่นที่มีรายละเอียดมากกว่าสัญญาณไซน์เพียงความถี่เดียว การวิเคราะห์รูปคลื่นเหล่านี้โดยพิจารณาเฉพาะค่า ณ เวลาใดเวลาหนึ่งอาจทำให้มองไม่เห็นองค์ประกอบความถี่ที่ซ่อนอยู่

แนวคิดของ \textbf{อนุกรมฟูเรียร์ (Fourier series)} ช่วยให้สามารถแทนฟังก์ชันเป็นคาบด้วยผลรวมขององค์ประกอบคงที่ สัญญาณโคไซน์ และสัญญาณไซน์ที่มีความถี่เป็นจำนวนเต็มเท่าของความถี่มูลฐาน แต่ละองค์ประกอบมีขนาดและเฟสเฉพาะตัว จึงเป็นสะพานเชื่อมระหว่างการมองสัญญาณใน \textbf{โดเมนเวลา (time domain)} กับการมองสัญญาณใน \textbf{โดเมนความถี่ (frequency domain)}

ก่อนเข้าสู่อนุกรมฟูเรียร์ ผู้เรียนจำเป็นต้องเข้าใจลำดับ อนุกรม ผลบวกย่อย และการลู่เข้า เพราะอนุกรมฟูเรียร์เป็นผลรวมอนันต์ชนิดหนึ่ง ความเข้าใจเรื่องการลู่เข้าทำให้สามารถตีความได้ว่า เมื่อเพิ่มจำนวนพจน์ฮาร์มอนิก รูปคลื่นประมาณจะเข้าใกล้รูปคลื่นต้นฉบับในความหมายใด และเหตุใดบริเวณจุดไม่ต่อเนื่องจึงยังเกิดการสั่นหรือการพุ่งเกินแม้เพิ่มจำนวนพจน์มากขึ้น

บทนี้จึงมุ่งสร้างฐานความรู้สำหรับบทที่ 5 ซึ่งจะนำอนุกรมฟูเรียร์ไปวิเคราะห์รูปคลื่นไฟฟ้า ค่าประสิทธิผล กำลังไฟฟ้า ฮาร์มอนิก และคุณภาพกำลังไฟฟ้าอย่างเป็นระบบ

%\par\noindent\rule{\textwidth}{0.4pt}\par

\subsection{2. ความรู้พื้นฐานก่อนเรียน}

ผู้เรียนควรมีความรู้และทักษะต่อไปนี้

\begin{enumerate}
\def\labelenumi{\arabic{enumi}.}
\item
  พีชคณิตพื้นฐาน การจัดรูปสมการ และการใช้สัญลักษณ์ผลรวม \(\sum\)
\item
  ฟังก์ชันตรีโกณมิติ \(\sin \theta\) และ \(\cos \theta\) รวมทั้งคาบ แอมพลิจูด และมุมเฟส
\item
  การอินทิเกรตฟังก์ชันพื้นฐานและการอินทิเกรตแบบจำกัดเขต
\item
  สูตรอินทิเกรตโดยส่วนและเอกลักษณ์ตรีโกณมิติพื้นฐาน
\item
  แนวคิดของฟังก์ชัน กราฟ ความต่อเนื่อง และจุดไม่ต่อเนื่อง
\item
  ความหมายของความถี่ \(f\) คาบ \(T\) และความถี่เชิงมุม \(\omega=2\pi f\)
\item
  ความรู้เบื้องต้นเกี่ยวกับแรงดัน กระแส และรูปคลื่นทางไฟฟ้า
\end{enumerate}

\begin{quote}
\textbf{ข้อเสนอแนะก่อนเริ่มบท:} ผู้สอนควรทบทวนอินทิกรัลของ \(\sin(n\omega t)\) และ \(\cos(n\omega t)\) รวมทั้งความสัมพันธ์ระหว่างคาบกับความถี่ประมาณ 15--20 นาที เพื่อให้ผู้เรียนพร้อมต่อการหาสัมประสิทธิ์ฟูเรียร์
\end{quote}

%\par\noindent\rule{\textwidth}{0.4pt}\par

\subsection{3. คำสำคัญประจำบท}

\begin{table}[htbp]
    \centering
    \begin{tabular}{lll}
        \hline
        \textbf{คำศัพท์ภาษาไทย} & \textbf{คำศัพท์ภาษาอังกฤษ} & \textbf{ความหมายโดยย่อ} \\
        \hline
        ลำดับ & Sequence & รายการของจำนวนหรือพจน์ที่จัดเรียงตามดัชนี \\
        พจน์ทั่วไป & General term & สูตรที่ใช้ระบุสมาชิกตำแหน่งที่ \(n\) ของลำดับ \\
        อนุกรม & Series & ผลบวกของพจน์ในลำดับ \\
        ผลบวกย่อย & Partial sum & ผลบวกของพจน์จำนวนจำกัดตั้งแต่พจน์แรกถึงพจน์ที่ \(N\) \\
        การลู่เข้า & Convergence & ภาวะที่ผลบวกย่อยเข้าใกล้ค่าจำกัดเมื่อจำนวนพจน์เพิ่มขึ้น \\
        การลู่ออก & Divergence & ภาวะที่ผลบวกย่อยไม่เข้าใกล้ค่าจำกัด \\
        ฟังก์ชันเป็นคาบ & Periodic function & ฟังก์ชันที่มีค่าเกิดซ้ำหลังเวลาหรือระยะคาบหนึ่ง \\
        คาบมูลฐาน & Fundamental period & คาบบวกที่น้อยที่สุดของฟังก์ชันเป็นคาบ \\
        ความถี่มูลฐาน & Fundamental frequency & ความถี่ต่ำสุดของสัญญาณเป็นคาบ มีค่า \(f_0=1/T\) \\
        ฮาร์มอนิก & Harmonic & องค์ประกอบไซน์หรือโคไซน์ที่มีความถี่เป็น \(n\) เท่าของความถี่มูลฐาน \\
        ออร์โทกอนัลลิตี & Orthogonality & สมบัติที่ผลคูณภายในของฟังก์ชันฐานต่างลำดับมีค่าเป็นศูนย์ \\
        สัมประสิทธิ์ฟูเรียร์ & Fourier coefficient & ค่าที่ระบุขนาดขององค์ประกอบฐานแต่ละความถี่ \\
        สเปกตรัมเส้น & Line spectrum & การแสดงขนาดหรือเฟสขององค์ประกอบที่ความถี่ไม่ต่อเนื่อง \\
        ฟังก์ชันคู่ & Even function & ฟังก์ชันที่มีสมบัติ \(f(-t)=f(t)\) \\
        ฟังก์ชันคี่ & Odd function & ฟังก์ชันที่มีสมบัติ \(f(-t)=-f(t)\) \\
        การสังเคราะห์ฟูเรียร์ & Fourier synthesis & การรวมองค์ประกอบฮาร์มอนิกเพื่อสร้างรูปคลื่นขึ้นใหม่ \\
        ปรากฏการณ์กิบส์ & Gibbs phenomenon & การพุ่งเกินและการสั่นใกล้จุดไม่ต่อเนื่องของผลรวมฟูเรียร์แบบตัดทอน \\
        \hline
    \end{tabular}
\end{table}

%\par\noindent\rule{\textwidth}{0.4pt}\par

\subsection{4. ผลลัพธ์การเรียนรู้ประจำบท}

\begin{quote}
เมื่อเรียนจบบทนี้ นักศึกษาสามารถ\ldots{}
\end{quote}

\begin{enumerate}
\def\labelenumi{\arabic{enumi}.}
\item
  \textbf{อธิบาย} ความหมายของลำดับ อนุกรม ผลบวกย่อย การลู่เข้า และการลู่ออกได้อย่างถูกต้อง
\item
  \textbf{คำนวณ} พจน์ทั่วไปและผลบวกของอนุกรมเลขคณิตและอนุกรมเรขาคณิตในกรณีที่กำหนดได้
\item
  \textbf{วิเคราะห์} คาบ ความถี่มูลฐาน และความถี่ฮาร์มอนิกของรูปคลื่นไฟฟ้าเป็นคาบได้
\item
  \textbf{คำนวณ} สัมประสิทธิ์ฟูเรียร์ตรีโกณมิติของฟังก์ชันเป็นคาบพื้นฐานโดยแสดงขั้นตอนอย่างเป็นระบบได้
\item
  \textbf{จำแนกและใช้} สมบัติของฟังก์ชันคู่และฟังก์ชันคี่เพื่อลดรูปการคำนวณอนุกรมฟูเรียร์ได้
\item
  \textbf{ตีความ} สเปกตรัมเส้น ฮาร์มอนิก และความหมายทางกายภาพขององค์ประกอบความถี่ในสัญญาณไฟฟ้าได้
\end{enumerate}

\subsubsection{4.1 ตารางความสอดคล้องของผลลัพธ์การเรียนรู้ กิจกรรม และการประเมิน}

\begin{table}[htbp]
    \centering
    \begin{tabular}{lllll}
        \hline
        \textbf{ผลลัพธ์การเรียนรู้ประจำบท} & \textbf{เนื้อหาที่เกี่ยวข้อง} & \textbf{กิจกรรมการเรียนรู้} & \textbf{วิธีประเมิน} & \textbf{หลักฐานการเรียนรู้} \\
        \hline
        CLO1 อธิบายลำดับ อนุกรม และการลู่เข้า & 7.1--7.3 & วิเคราะห์ผลบวกย่อยจากตารางและกราฟ & คำถามสั้นและแบบฝึกหัดพื้นฐาน & ใบงานนิยามและกราฟผลบวกย่อย \\
        CLO2 คำนวณอนุกรมเลขคณิตและเรขาคณิต & 7.2--7.3 & คำนวณร่วมกันและตรวจคำตอบแบบเพื่อนช่วยเพื่อน & แบบฝึกหัดคำนวณ & วิธีทำที่แสดงพจน์ทั่วไปและผลบวก \\
        CLO3 วิเคราะห์คาบและฮาร์มอนิก & 7.4, 7.10 & อ่านกราฟรูปคลื่นและสร้างตารางความถี่ฮาร์มอนิก & แบบทดสอบระหว่างเรียน & ตาราง \(T\), \(f_0\), \(\omega_0\), \(nf_0\) \\
        CLO4 คำนวณสัมประสิทธิ์ฟูเรียร์ & 7.5--7.8, 7.13 & แก้โจทย์แบบมีโครงขั้นตอน & แบบฝึกหัดวิเคราะห์และข้อสอบคำนวณ & สมการอินทิกรัลและคำตอบสัมประสิทธิ์ \\
        CLO5 ใช้สมบัติฟังก์ชันคู่และคี่ & 7.9 & จำแนกกราฟและลดรูปอินทิกรัล & ใบงานวิเคราะห์สมมาตร & ตารางการตัดสัมประสิทธิ์ที่เป็นศูนย์ \\
        CLO6 ตีความสเปกตรัมและความหมายทางกายภาพ & 7.10--7.12, 9 & วิเคราะห์กรณีรูปคลื่นไฟฟ้าและอภิปราย & งานกลุ่มและคำถามปลายเปิด & แผนภาพสเปกตรัมพร้อมคำอธิบาย \\
        \hline
    \end{tabular}
\end{table}

%\par\noindent\rule{\textwidth}{0.4pt}\par

\subsection{5. แผนผังเนื้อหาของบท}

ลำดับการเรียนรู้ของบทนี้สามารถสรุปได้ดังนี้

\textbf{ลำดับ} \(\rightarrow\) \textbf{อนุกรมและผลบวกย่อย} \(\rightarrow\) \textbf{การลู่เข้า} \(\rightarrow\) \textbf{ฟังก์ชันเป็นคาบ} \(\rightarrow\) \textbf{ฐานไซน์--โคไซน์และออร์โทกอนัลลิตี} \(\rightarrow\) \textbf{สัมประสิทธิ์ฟูเรียร์} \(\rightarrow\) \textbf{สมมาตรคู่--คี่} \(\rightarrow\) \textbf{ฮาร์มอนิกและสเปกตรัม} \(\rightarrow\) \textbf{การสังเคราะห์รูปคลื่น} \(\rightarrow\) \textbf{การตีความเชิงวิศวกรรมไฟฟ้า}

% [รูปที่ 1: แผนผังความเชื่อมโยงจากอนุกรมสู่การวิเคราะห์สเปกตรัม]
% ประเภทภาพ: Flowchart / Infographic
% ตำแหน่งที่แนะนำ: หลังหัวข้อ 5 แผนผังเนื้อหาของบท
% วัตถุประสงค์การเรียนรู้: ช่วยให้ผู้เรียนมองเห็นว่าความรู้เรื่องอนุกรมและการลู่เข้าเป็นฐานของอนุกรมฟูเรียร์และการวิเคราะห์ความถี่
% คำอธิบายภาพ: แสดงกล่องเรียงจากซ้ายไปขวา ได้แก่ Sequence, Series, Partial Sum, Convergence, Periodic Function, Orthogonal Sinusoids, Fourier Coefficients, Harmonic Spectrum และ Electrical Waveform Analysis พร้อมลูกศรแสดงลำดับเหตุผล
% องค์ประกอบที่ต้องมี: สัญลักษณ์ $a_n$, $S_N$, $T$, $f_0$, $a_n,b_n$, แท่งสเปกตรัม และรูปคลื่นสี่เหลี่ยม
% ข้อความหรือป้ายกำกับ: “ลำดับ”, “ผลบวกย่อย”, “การลู่เข้า”, “คาบ”, “สัมประสิทธิ์”, “ฮาร์มอนิก”, “โดเมนเวลา”, “โดเมนความถี่”
% ภาษาของป้ายกำกับ: ไทยและอังกฤษ
% สัดส่วนภาพ: 16:9
% Image Generation Prompt: Clean academic engineering flowchart showing the learning progression from sequence and series to Fourier analysis: Sequence a_n → Series and partial sum S_N → Convergence → Periodic function with period T → Orthogonal sine and cosine basis → Fourier coefficients a_n and b_n → Harmonic line spectrum → Electrical waveform analysis. Include a small square-wave icon in time domain and discrete spectral bars in frequency domain. Thai-English labels, white background, blue technical textbook style, clear arrows, no decorative clutter, 16:9.
% Alt Text: แผนผังแสดงลำดับความรู้จากลำดับและผลบวกย่อยไปสู่อนุกรมฟูเรียร์ สเปกตรัมฮาร์มอนิก และการวิเคราะห์รูปคลื่นไฟฟ้า

\begin{figure}[htbp]
    \centering
    \includegraphics[width=0.8\textwidth]{images/figure-01.png}
    \caption{แผนผังความเชื่อมโยงจากอนุกรมสู่การวิเคราะห์สเปกตรัม}
\end{figure}

%\par\noindent\rule{\textwidth}{0.4pt}\par

\subsection{6. บทนำ}

ในระบบไฟฟ้ากระแสสลับอุดมคติ แรงดันและกระแสอาจอธิบายด้วยฟังก์ชันไซน์เพียงความถี่เดียว อย่างไรก็ตาม อุปกรณ์จริงจำนวนมากทำงานด้วยการเปิด--ปิดสวิตช์ การเรียงกระแส หรือการควบคุมด้วยอุปกรณ์สารกึ่งตัวนำ จึงทำให้รูปคลื่นเบี่ยงเบนจากไซน์ ตัวอย่างเช่น แรงดันเอาต์พุตของอินเวอร์เตอร์แบบพัลส์ กระแสอินพุตของแหล่งจ่ายไฟสวิตชิ่ง และสัญญาณนาฬิกาดิจิทัล ล้วนมีขอบชัน จุดไม่ต่อเนื่อง หรือช่วงค่าคงที่

การพิจารณารูปคลื่นเฉพาะในโดเมนเวลาช่วยให้ทราบว่าแรงดันหรือกระแสเปลี่ยนแปลงอย่างไรตามเวลา แต่ยังไม่ตอบโดยตรงว่ารูปคลื่นนั้นประกอบด้วยความถี่ใดบ้าง อนุกรมฟูเรียร์ให้วิธีแยกรูปคลื่นเป็นผลรวมขององค์ประกอบที่เป็นไซน์และโคไซน์ ซึ่งแต่ละองค์ประกอบสามารถนำไปวิเคราะห์ผ่านวงจรหรือระบบเชิงเส้นได้ง่ายกว่า จากนั้นจึงรวมผลตอบสนองกลับด้วยหลักการซ้อนทับ

แนวคิดนี้มีความสำคัญต่อการวิเคราะห์สัญญาณ ระบบสื่อสาร ระบบควบคุม อิเล็กทรอนิกส์กำลัง คุณภาพกำลังไฟฟ้า และการประมวลผลสัญญาณ รายวิชานี้วางบทที่ 4 เป็นพื้นฐานทางคณิตศาสตร์ ก่อนขยายสู่การประยุกต์วิเคราะห์รูปคลื่นและฮาร์มอนิกอย่างละเอียดในบทที่ 5

MIT OpenCourseWare อธิบายอนุกรมฟูเรียร์ในฐานะการแทนฟังก์ชันเป็นคาบด้วยองค์ประกอบไซน์และโคไซน์ที่สัมพันธ์กันทางฮาร์มอนิก ส่วนเอกสารด้านสัญญาณและระบบของ Stanford และ Rice University เชื่อมโยงสัมประสิทธิ์ฟูเรียร์กับการมองสัญญาณในโดเมนความถี่ การตีความเช่นนี้สอดคล้องกับงานวิศวกรรมไฟฟ้าที่ต้องพิจารณาทั้งรูปร่างคลื่นและองค์ประกอบความถี่ (MIT OpenCourseWare, 2011; Osgood, 2007; Johnson, 2016)

%\par\noindent\rule{\textwidth}{0.4pt}\par

\subsection{7. เนื้อหาหลัก}

\subsubsection{7.1 ความหมายของลำดับและอนุกรม}

\paragraph{7.1.1 ลำดับ}

\textbf{ลำดับ (sequence)} คือฟังก์ชันที่กำหนดค่าตามดัชนีจำนวนเต็มบวกหรือจำนวนเต็มไม่ลบ โดยมักเขียนเป็น

\[
\{a_n\}=a_1,a_2,a_3,\ldots
\]

หรือหากเริ่มที่ดัชนีศูนย์ เขียนเป็น

\[
\{a_n\}_{n=0}^{\infty}=a_0,a_1,a_2,\ldots
\]

ตัวแปร \(n\) เป็นดัชนี ไม่มีหน่วย ส่วน \(a_n\) อาจมีหน่วยตามปริมาณที่แทน เช่น โวลต์ แอมแปร์ หรือวัตต์

\textbf{พจน์ทั่วไป (general term)} คือสมการที่ใช้คำนวณพจน์ตำแหน่งใด ๆ เช่น

\[
a_n=2n+1
\]

ทำให้ได้ลำดับ \(3,5,7,9,\ldots\)

ในงานวิศวกรรม ลำดับอาจเกิดจากค่าตัวอย่างของสัญญาณตามเวลา ค่าฮาร์มอนิกลำดับที่ \(n\) หรือค่าประมาณที่ได้จากการคำนวณซ้ำ ดังนั้นดัชนี \(n\) มิได้หมายถึงเวลาเสมอไป แต่เป็นตัวระบุตำแหน่งของข้อมูลหรือองค์ประกอบ

\paragraph{7.1.2 อนุกรม}

\textbf{อนุกรม (series)} คือผลบวกของพจน์ในลำดับ เช่น

\[
\sum_{n=1}^{\infty}a_n=a_1+a_2+a_3+\cdots
\]

เครื่องหมาย \(\infty\) มิได้หมายความว่าสามารถบวกครบทุกพจน์จริง แต่เป็นการศึกษาพฤติกรรมของผลบวกเมื่อเพิ่มจำนวนพจน์มากขึ้น แนวคิดสำคัญจึงอยู่ที่ \textbf{ผลบวกย่อย (partial sum)}

\[
S_N=\sum_{n=1}^{N}a_n
\]

โดย \(S_N\) เป็นผลบวกของ \(N\) พจน์แรก หาก \(S_N\) เข้าใกล้ค่าคงที่ \(S\) เมื่อ \(N\) เพิ่มโดยไม่สิ้นสุด จะกล่าวว่าอนุกรมลู่เข้าสู่ \(S\)

\[
\lim_{N\to\infty}S_N=S
\]

\textbf{หน่วย:} ถ้า \(a_n\) มีหน่วยโวลต์ ผลบวกย่อย \(S_N\) ก็มีหน่วยโวลต์ แต่สำหรับอนุกรมฟูเรียร์แต่ละพจน์เป็นค่าฟังก์ชัน ณ เวลาเดียวกัน จึงรวมกันเพื่อให้ได้หน่วยเดียวกับสัญญาณต้นฉบับ

% [รูปที่ 2: ความแตกต่างระหว่างลำดับ อนุกรม และผลบวกย่อย]
% ประเภทภาพ: Infographic / Graph
% ตำแหน่งที่แนะนำ: หลังหัวข้อ 7.1.2
% วัตถุประสงค์การเรียนรู้: แยกความหมายของสมาชิก $a_n$ ออกจากผลบวกย่อย $S_N$
% คำอธิบายภาพ: ด้านซ้ายเป็นแท่งแสดงพจน์ $a_1,a_2,\ldots$ ด้านขวาเป็นกราฟ $S_N$ เทียบกับ $N$ ที่เข้าใกล้เส้นค่าจำกัด $S$
% องค์ประกอบที่ต้องมี: แกนนอน $n$ หรือ $N$ แกนตั้ง $a_n$ และ $S_N$ เส้นประ $S$ และลูกศร “เพิ่มจำนวนพจน์”
% ข้อความหรือป้ายกำกับ: “สมาชิกของลำดับ”, “ผลบวกย่อย”, “ค่าลิมิต”
% ภาษาของป้ายกำกับ: ไทยและอังกฤษ
% สัดส่วนภาพ: 16:9
% Image Generation Prompt: Educational two-panel infographic. Left panel: a sequence shown as individual vertical bars labeled a1, a2, a3, ... . Right panel: a partial-sum graph S_N versus N approaching a horizontal dashed limit S. Clearly distinguish individual terms from cumulative sums. Thai-English labels, clean engineering mathematics textbook style, white background, 16:9.
% Graph Generation Prompt: Plot an example sequence $a_n=(1/2)^n$ for $n=0$ to $10$ and the corresponding partial sums $S_N=\sum_{n=0}^{N}(1/2)^n$ for $N=0$ to $10$. Use horizontal axis “index n or N”, vertical axes “term a_n” and “partial sum S_N”, and a dashed line at $S=2$.
% Alt Text: ภาพเปรียบเทียบพจน์ของลำดับที่ลดลงกับผลบวกย่อยที่เพิ่มขึ้นและเข้าใกล้ค่าจำกัด

\begin{figure}[htbp]
    \centering
    \includegraphics[width=0.8\textwidth]{images/figure-02.png}
    \caption{ความแตกต่างระหว่างลำดับ อนุกรม และผลบวกย่อย}
\end{figure}

\subsubsection{7.2 อนุกรมเลขคณิตและอนุกรมเรขาคณิต}

\paragraph{7.2.1 ลำดับและอนุกรมเลขคณิต}

ลำดับเลขคณิตมีผลต่างร่วม \(d\) คงที่

\[
a_n=a_1+(n-1)d
\]

ผลบวก \(N\) พจน์แรกคือ

\[
S_N=\frac{N}{2}\left[2a_1+(N-1)d\right]
\]

หรือ

\[
S_N=\frac{N}{2}(a_1+a_N)
\]

\textbf{ที่มาโดยย่อ:} เขียนผลบวกไปข้างหน้าและย้อนกลับ แล้วบวกสมการเป็นคู่ แต่ละคู่มีค่า \(a_1+a_N\) เท่ากันจำนวน \(N\) คู่ จากนั้นหารสอง

\textbf{เงื่อนไขการใช้:} ต้องเป็นลำดับที่ผลต่างระหว่างพจน์ติดกันคงที่

\textbf{ข้อจำกัด:} อนุกรมเลขคณิตอนันต์ที่พจน์ไม่เป็นศูนย์ทั้งหมดโดยทั่วไปไม่ลู่เข้า เพราะพจน์ \(a_n\) ไม่เข้าใกล้ศูนย์หรือผลบวกย่อยเติบโตโดยไม่มีขอบเขต

\paragraph{7.2.2 ลำดับและอนุกรมเรขาคณิต}

ลำดับเรขาคณิตมีอัตราส่วนร่วม \(r\) คงที่

\[
a_n=a_1r^{n-1}
\]

ผลบวก \(N\) พจน์แรก เมื่อ \(r\ne1\) คือ

\[
S_N=a_1\frac{1-r^N}{1-r}
\]

หากเริ่มลำดับที่ \(n=0\) ด้วยพจน์ \(a_0\) จะเขียนได้ว่า

\[
\sum_{n=0}^{N}a_0r^n=a_0\frac{1-r^{N+1}}{1-r}
\]

สำหรับอนุกรมอนันต์ หาก \(|r|<1\) จะมี \(r^N\to0\) และได้

\[
\sum_{n=1}^{\infty}a_1r^{n-1}=\frac{a_1}{1-r}
\]

\textbf{ความหมายทางวิศวกรรม:} รูปแบบเรขาคณิตเกิดในกระบวนการที่ลดลงด้วยสัดส่วนคงที่ เช่น ค่าประมาณจากการสะท้อนซ้ำ การลดทอนต่อช่วง หรือความคลาดเคลื่อนของวิธีคำนวณซ้ำบางชนิด

\paragraph{ตัวอย่างที่ 1: ผลบวกของอนุกรมเลขคณิต}

\textbf{ตัวอย่างเพิ่มเติมเพื่อการเรียนรู้}\\
ลำดับ \(5,8,11,14,\ldots\) จงหาพจน์ที่ 20 และผลบวก 20 พจน์แรก

\textbf{1. วิเคราะห์โจทย์}\\
เป็นลำดับเลขคณิต มี \(a_1=5\), \(d=3\), \(N=20\)

\textbf{2. หาพจน์ที่ 20}

\[
a_{20}=a_1+(20-1)d=5+19(3)=62
\]

\textbf{3. หาผลบวก}

\[
S_{20}=\frac{20}{2}(5+62)=10(67)=670
\]

\textbf{4. ตรวจสอบความสมเหตุสมผล}\\
ค่าเฉลี่ยของพจน์แรกและพจน์สุดท้ายคือ \(33.5\) เมื่อคูณ 20 ได้ \(670\) สอดคล้องกัน

\textbf{คำตอบ:} \(a_{20}=62\) และ \(S_{20}=670\)

\paragraph{ตัวอย่างที่ 2: อนุกรมเรขาคณิตของค่าที่ลดลง}

\textbf{ตัวอย่างเพิ่มเติมเพื่อการเรียนรู้}\\
ค่าความคลาดเคลื่อนของการประมาณหนึ่งมีขนาด \(10,5,2.5,1.25,\ldots\) จงหาผลรวมความคลาดเคลื่อนตามแบบจำลองอนุกรมอนันต์

\textbf{1. กำหนดตัวแปร}\\
\(a_1=10\) และ \(r=0.5\)

\textbf{2. ตรวจสอบการลู่เข้า}

\[
|r|=0.5<1
\]

จึงลู่เข้า

\textbf{3. คำนวณผลบวก}

\[
S=\frac{a_1}{1-r}=\frac{10}{1-0.5}=20
\]

\textbf{4. แปลความหมาย}\\
แม้มีจำนวนพจน์ไม่สิ้นสุด แต่ผลรวมมีค่าจำกัดเท่ากับ 20 หน่วย เพราะแต่ละพจน์ลดลงอย่างรวดเร็ว

\subsubsection{7.3 แนวคิดของการลู่เข้าและการลู่ออกของอนุกรม}

\paragraph{7.3.1 เงื่อนไขจำเป็นเบื้องต้น}

ถ้าอนุกรม \(\sum a_n\) ลู่เข้า จำเป็นต้องมี

\[
\lim_{n\to\infty}a_n=0
\]

แต่เงื่อนไขนี้ \textbf{ยังไม่เพียงพอ} ตัวอย่างเช่น อนุกรมฮาร์มอนิก

\[
\sum_{n=1}^{\infty}\frac{1}{n}
\]

มีพจน์ \(1/n\to0\) แต่ผลบวกย่อยยังเพิ่มขึ้นโดยไม่เข้าใกล้ค่าจำกัด จึงลู่ออก

ข้อสังเกตนี้สำคัญมาก เพราะผู้เรียนมักสรุปผิดว่า ``พจน์เข้าใกล้ศูนย์จึงอนุกรมลู่เข้า'' ความจริงคือ ถ้าพจน์ไม่เข้าใกล้ศูนย์ยืนยันได้ว่าอนุกรมลู่ออก แต่ถ้าพจน์เข้าใกล้ศูนย์ยังต้องตรวจสอบเพิ่มเติม

\paragraph{7.3.2 การลู่เข้าของอนุกรมเรขาคณิต}

จาก

\[
S_N=a_1\frac{1-r^N}{1-r}
\]

พฤติกรรมขึ้นกับ \(r\)

\begin{itemize}
\item
  ถ้า \(|r|<1\) แล้ว \(r^N\to0\) อนุกรมลู่เข้า
\item
  ถ้า \(r=1\) ผลบวกเพิ่มเป็น \(Na_1\) จึงลู่ออกเมื่อ \(a_1\ne0\)
\item
  ถ้า \(r=-1\) ผลบวกสลับไปมาและไม่เข้าใกล้ค่าหนึ่ง
\item
  ถ้า \(|r|>1\) ขนาดพจน์เพิ่มขึ้นและอนุกรมลู่ออก
\end{itemize}

\paragraph{7.3.3 ความหมายของการลู่เข้าสำหรับอนุกรมฟังก์ชัน}

สำหรับอนุกรมของฟังก์ชัน เช่น อนุกรมฟูเรียร์ เราสนใจว่า

\[
S_N(t)=\frac{a_0}{2}+\sum_{n=1}^{N}\left[a_n\cos(n\omega_0t)+b_n\sin(n\omega_0t)\right]
\]

เข้าใกล้ \(f(t)\) อย่างไรเมื่อ \(N\) เพิ่มขึ้น การลู่เข้าอาจพิจารณาได้หลายแบบ เช่น

\begin{enumerate}
\def\labelenumi{\arabic{enumi}.}
\item
  \textbf{การลู่เข้าทีละจุด (pointwise convergence):} สำหรับแต่ละค่า \(t\) ผลบวกย่อยเข้าใกล้ค่าหนึ่ง
\item
  \textbf{การลู่เข้าแบบสม่ำเสมอ (uniform convergence):} ค่าคลาดเคลื่อนสูงสุดตลอดช่วงเข้าใกล้ศูนย์
\item
  \textbf{การลู่เข้าในความหมายกำลังสองเฉลี่ย (mean-square convergence):} ค่าเฉลี่ยของกำลังสองความคลาดเคลื่อนเข้าใกล้ศูนย์
\end{enumerate}

ในบทนี้เน้นการตีความเชิงวิศวกรรมว่า การเพิ่มจำนวนฮาร์มอนิกทำให้รูปคลื่นประมาณใกล้เคียงต้นฉบับมากขึ้นโดยรวม แต่ที่จุดกระโดดอาจไม่ลู่เข้าหาค่าด้านซ้ายหรือด้านขวาโดยตรง และอาจเกิดการพุ่งเกินบริเวณใกล้จุดนั้น

% [รูปที่ 3: การลู่เข้าของผลบวกย่อยและตัวอย่างการลู่ออก]
% ประเภทภาพ: Graph
% ตำแหน่งที่แนะนำ: หลังหัวข้อ 7.3.3
% วัตถุประสงค์การเรียนรู้: เปรียบเทียบพฤติกรรมของผลบวกย่อยที่เข้าใกล้ค่าจำกัดกับผลบวกย่อยที่เพิ่มไม่สิ้นสุดหรือสลับไปมา
% คำอธิบายภาพ: กราฟสามช่อง แสดง $S_N$ ของอนุกรมเรขาคณิต $r=0.5$ ที่เข้าใกล้ 2, อนุกรมฮาร์มอนิกที่เพิ่มช้าแต่ไม่จำกัด และอนุกรม $1-1+1-1+\cdots$ ที่สลับระหว่าง 0 กับ 1
% องค์ประกอบที่ต้องมี: แกนนอน $N$ แกนตั้ง $S_N$ เส้นค่าจำกัด และคำกำกับ Convergent/Divergent/Oscillatory
% ข้อความหรือป้ายกำกับ: “ลู่เข้า”, “ลู่ออก”, “สลับไม่ลู่เข้า”
% ภาษาของป้ายกำกับ: ไทยและอังกฤษ
% สัดส่วนภาพ: 16:9
% Image Generation Prompt: Three-panel mathematical graph comparing partial sums. Panel 1: geometric series with r=0.5 approaching S=2. Panel 2: harmonic series partial sums increasing slowly without a finite limit. Panel 3: alternating sequence of partial sums 1,0,1,0 showing oscillatory divergence. Label axes N and S_N, Thai-English annotations, clean textbook style, white background, 16:9.
% Graph Generation Prompt: Plot partial sums for $\sum_{n=0}^{N}(1/2)^n$, $\sum_{n=1}^{N}1/n$, and $\sum_{n=0}^{N}(-1)^n$ over $N=1$ to $50$. Use separate panels, axis labels, legends, and a dashed line $S=2$ in the first panel.
% Alt Text: กราฟเปรียบเทียบผลบวกย่อยของอนุกรมที่ลู่เข้า อนุกรมที่เพิ่มไม่สิ้นสุด และอนุกรมที่สลับไปมา

\begin{figure}[htbp]
    \centering
    \includegraphics[width=0.8\textwidth]{images/figure-03.png}
    \caption{การลู่เข้าของผลบวกย่อยและตัวอย่างการลู่ออก}
\end{figure}

\subsubsection{7.4 ฟังก์ชันเป็นคาบและรูปคลื่นทางไฟฟ้า}

\paragraph{7.4.1 นิยามของฟังก์ชันเป็นคาบ}

ฟังก์ชัน \(f(t)\) เป็นฟังก์ชันเป็นคาบ หากมีจำนวนบวก \(T\) ที่ทำให้

\[
f(t+T)=f(t)
\]

สำหรับทุกค่า \(t\) ที่อยู่ในโดเมน ค่าบวกที่น้อยที่สุดซึ่งทำให้สมการเป็นจริงเรียกว่า \textbf{คาบมูลฐาน (fundamental period)}

ความถี่มูลฐานคือ

\[
f_0=\frac{1}{T}
\]

และความถี่เชิงมุมมูลฐานคือ

\[
\omega_0=2\pi f_0=\frac{2\pi}{T}
\]

\textbf{หน่วย:}

\begin{itemize}
\item
  \(T\) มีหน่วยวินาที (s)
\item
  \(f_0\) มีหน่วยเฮิรตซ์ (Hz)
\item
  \(\omega_0\) มีหน่วยเรเดียนต่อวินาที (rad/s)
\end{itemize}

\paragraph{7.4.2 ตัวอย่างรูปคลื่นเป็นคาบ}

รูปคลื่นเป็นคาบที่พบบ่อย ได้แก่

\begin{itemize}
\item
  รูปคลื่นไซน์และโคไซน์
\item
  รูปคลื่นสี่เหลี่ยม
\item
  รูปคลื่นสามเหลี่ยม
\item
  รูปคลื่นฟันเลื่อย
\item
  ขบวนพัลส์
\item
  รูปคลื่นแรงดันหรือกระแสหลังการเรียงกระแส
\item
  รูปคลื่น PWM ที่มีรูปแบบซ้ำอย่างคงที่
\end{itemize}

รูปคลื่นบางชนิดมีคาบจากการรวมหลายสัญญาณ หากอัตราส่วนความถี่ขององค์ประกอบเป็นจำนวนตรรกยะ สัญญาณรวมจะมีคาบร่วม แต่ถ้าอัตราส่วนความถี่เป็นจำนวนอตรรกยะ สัญญาณรวมโดยทั่วไปจะไม่เป็นคาบอย่างแท้จริง

\paragraph{7.4.3 ฮาร์มอนิกของฟังก์ชันเป็นคาบ}

เมื่อสัญญาณมีคาบ \(T\) ความถี่ที่ปรากฏในอนุกรมฟูเรียร์จะอยู่ที่

\[
f_n=nf_0,\qquad n=1,2,3,\ldots
\]

หรือในรูปความถี่เชิงมุม

\[
\omega_n=n\omega_0
\]

องค์ประกอบที่ \(n=1\) เรียกว่า \textbf{ฮาร์มอนิกมูลฐาน} ส่วน \(n=2,3,\ldots\) เรียกว่า \textbf{ฮาร์มอนิกลำดับที่ 2, 3, \ldots{}} องค์ประกอบคงที่หรือค่าเฉลี่ยมักเรียกว่าองค์ประกอบ DC และอยู่ที่ความถี่ศูนย์

\paragraph{ตัวอย่างที่ 3: การหาความถี่มูลฐานและฮาร์มอนิก}

\textbf{ตัวอย่างเพิ่มเติมเพื่อการเรียนรู้}\\
รูปคลื่นแรงดันหนึ่งมีคาบ \(T=20\ \text{ms}\) จงหา \(f_0\), \(\omega_0\) และความถี่ฮาร์มอนิกลำดับที่ 3 และ 5

\textbf{1. แปลงหน่วย}

\[
T=20\times10^{-3}=0.020\ \text{s}
\]

\textbf{2. ความถี่มูลฐาน}

\[
f_0=\frac{1}{0.020}=50\ \text{Hz}
\]

\textbf{3. ความถี่เชิงมุมมูลฐาน}

\[
\omega_0=2\pi(50)=100\pi\approx314.16\ \text{rad/s}
\]

\textbf{4. ฮาร์มอนิก}

\[
f_3=3f_0=150\ \text{Hz}
\]

\[
f_5=5f_0=250\ \text{Hz}
\]

\textbf{การแปลความหมาย:} หากรูปคลื่นนี้มีฮาร์มอนิกลำดับที่ 3 และ 5 องค์ประกอบดังกล่าวจะปรากฏที่ 150 Hz และ 250 Hz ตามลำดับ ไม่ใช่ความถี่สุ่มใด ๆ

% [รูปที่ 4: รูปคลื่นเป็นคาบและความสัมพันธ์ระหว่างคาบกับความถี่]
% ประเภทภาพ: Graph / Technical Illustration
% ตำแหน่งที่แนะนำ: หลังตัวอย่างที่ 3
% วัตถุประสงค์การเรียนรู้: เชื่อมโยงการเกิดซ้ำในโดเมนเวลากับความถี่มูลฐานและฮาร์มอนิก
% คำอธิบายภาพ: แสดงรูปคลื่นสี่เหลี่ยมและรูปคลื่นสามเหลี่ยมในโดเมนเวลา ทำเครื่องหมายช่วงคาบ $T$ พร้อมแผงข้างแสดงเส้นสเปกตรัมที่ $f_0,2f_0,3f_0$
% องค์ประกอบที่ต้องมี: แกนเวลา $t$ แกนแอมพลิจูด ช่วง $T$ ลูกศร $f_0=1/T$ และแท่งฮาร์มอนิก
% ข้อความหรือป้ายกำกับ: “คาบมูลฐาน”, “ความถี่มูลฐาน”, “ฮาร์มอนิก”
% ภาษาของป้ายกำกับ: ไทยและอังกฤษ
% สัดส่วนภาพ: 16:9
% Image Generation Prompt: Engineering textbook illustration linking time and frequency domains. Left: periodic square and triangular waveforms with one fundamental period T clearly marked. Right: discrete line spectrum at f0, 2f0, 3f0, and higher harmonics. Include equation f0=1/T and arrows connecting repetition in time to harmonic frequencies. Thai-English labels, white background, clean vector style, 16:9.
% Graph Generation Prompt: Draw a periodic square wave of period $T=20$ ms over three cycles and a line spectrum at 50, 100, 150, 200, and 250 Hz. Label the time axis in milliseconds, frequency axis in hertz, and mark the fundamental and harmonic orders.
% Alt Text: รูปคลื่นเกิดซ้ำทุกคาบ T เชื่อมกับสเปกตรัมเส้นที่ความถี่เป็นจำนวนเต็มเท่าของความถี่มูลฐาน

\begin{figure}[htbp]
    \centering
    \includegraphics[width=0.8\textwidth]{images/figure-04.png}
    \caption{รูปคลื่นเป็นคาบและความสัมพันธ์ระหว่างคาบกับความถี่}
\end{figure}

\subsubsection{7.5 แนวคิดพื้นฐานของอนุกรมฟูเรียร์}

\paragraph{7.5.1 การแทนรูปคลื่นด้วยผลรวมของไซน์และโคไซน์}

อนุกรมฟูเรียร์ตั้งอยู่บนแนวคิดว่า ฟังก์ชันเป็นคาบที่มีความเหมาะสมทางคณิตศาสตร์สามารถแทนหรือประมาณด้วยผลรวมของฟังก์ชันไซน์และโคไซน์ที่มีความถี่สัมพันธ์กันเป็นจำนวนเต็มเท่าของความถี่มูลฐานได้

สำหรับฟังก์ชันจริง \(f(t)\) ที่มีคาบ \(T\) รูปมาตรฐานของอนุกรมฟูเรียร์ตรีโกณมิติคือ

\[
f(t)=\frac{a_0}{2}+\sum_{n=1}^{\infty}\left[a_n\cos(n\omega_0t)+b_n\sin(n\omega_0t)\right]
\]

โดย

\begin{itemize}
\item
  \(a_0/2\) คือองค์ประกอบค่าเฉลี่ยหรือ DC
\item
  \(a_n\) คือสัมประสิทธิ์ขององค์ประกอบโคไซน์ลำดับที่ \(n\)
\item
  \(b_n\) คือสัมประสิทธิ์ขององค์ประกอบไซน์ลำดับที่ \(n\)
\item
  \(n\omega_0\) คือความถี่เชิงมุมของฮาร์มอนิกลำดับที่ \(n\)
\item
  \(n=1,2,3,\ldots\)
\end{itemize}

สัมประสิทธิ์ทุกตัวมีหน่วยเดียวกับ \(f(t)\) เช่น ถ้า \(f(t)\) เป็นแรงดัน สัมประสิทธิ์มีหน่วยโวลต์

อนุกรมฟูเรียร์มิได้กล่าวว่ารูปคลื่นไม่เป็นไซน์ ``กลายเป็น'' ไซน์ แต่กล่าวว่าเราสามารถสร้างรูปคลื่นนั้นขึ้นจากการซ้อนทับองค์ประกอบไซน์และโคไซน์จำนวนมาก เมื่อรวมองค์ประกอบอย่างเหมาะสม รูปคลื่นผลรวมจะเข้าใกล้รูปคลื่นต้นฉบับ

\paragraph{7.5.2 การวิเคราะห์และการสังเคราะห์}

กระบวนการฟูเรียร์ประกอบด้วยสองมุมมอง

\begin{enumerate}
\def\labelenumi{\arabic{enumi}.}
\item
  \textbf{การวิเคราะห์ฟูเรียร์ (Fourier analysis):} เริ่มจาก \(f(t)\) แล้วหาค่า \(a_0\), \(a_n\), \(b_n\)
\item
  \textbf{การสังเคราะห์ฟูเรียร์ (Fourier synthesis):} นำสัมประสิทธิ์ที่ได้มารวมเป็นผลบวกเพื่อสร้าง \(f(t)\) หรือค่าประมาณ \(S_N(t)\)
\end{enumerate}

หากใช้จำนวนฮาร์มอนิกจำกัด \(N\) จะได้ผลรวมแบบตัดทอน

\[
S_N(t)=\frac{a_0}{2}+\sum_{n=1}^{N}\left[a_n\cos(n\omega_0t)+b_n\sin(n\omega_0t)\right]
\]

ในงานจริง เราไม่สามารถรวมพจน์อนันต์ จึงเลือกจำนวนพจน์ตามความแม่นยำที่ต้องการ แบนด์วิดท์ของระบบ และข้อจำกัดในการคำนวณ

% [รูปที่ 5: การสังเคราะห์รูปคลื่นสี่เหลี่ยมจากฮาร์มอนิก]
% ประเภทภาพ: Graph / Infographic
% ตำแหน่งที่แนะนำ: หลังหัวข้อ 7.5.2
% วัตถุประสงค์การเรียนรู้: แสดงให้เห็นว่าการเพิ่มฮาร์มอนิกคี่ทำให้ผลรวมเข้าใกล้รูปคลื่นสี่เหลี่ยม
% คำอธิบายภาพ: แสดงกราฟ 4 แผง ได้แก่ ฮาร์มอนิกมูลฐานเพียงพจน์เดียว ผลรวม 3 พจน์ ผลรวม 5 พจน์ และผลรวม 15 พจน์ เปรียบเทียบกับเส้นรูปคลื่นสี่เหลี่ยมอุดมคติ
% องค์ประกอบที่ต้องมี: สมการ $S_N(t)=\frac{4V_m}{\pi}\sum_{k=0}^{N-1}\frac{1}{2k+1}\sin[(2k+1)\omega_0t]$ แกนเวลา แอมพลิจูด $\pm V_m$ และจุดกระโดด
% ข้อความหรือป้ายกำกับ: “1 harmonic”, “3 harmonics”, “5 harmonics”, “15 harmonics”, “ideal square wave”
% ภาษาของป้ายกำกับ: อังกฤษ หรือไทยและอังกฤษ
% สัดส่วนภาพ: 16:9
% Image Generation Prompt: Four-panel engineering mathematics graph showing Fourier synthesis of an odd square wave. Compare the ideal square wave with partial sums using 1, 3, 5, and 15 odd harmonics. Use consistent axes, mark amplitude +Vm and -Vm, show increasingly sharp transitions and small ringing near discontinuities. White background, clean textbook graph, 16:9.
% Graph Generation Prompt: Plot $S_N(t)=\frac{4V_m}{\pi}\sum_{k=0}^{N-1}\frac{1}{2k+1}\sin[(2k+1)\omega_0t]$ for $N=1,3,5,15$, with $V_m=1$, $T=2\pi$, and $t\in[-\pi,\pi]$. Overlay the ideal odd square wave. Label axes and legends.
% Alt Text: ผลรวมของฮาร์มอนิกคี่จำนวนเพิ่มขึ้นทำให้รูปคลื่นประมาณเปลี่ยนจากไซน์ไปใกล้รูปคลื่นสี่เหลี่ยม

\begin{figure}[htbp]
    \centering
    \includegraphics[width=0.8\textwidth]{images/figure-05.png}
    \caption{การสังเคราะห์รูปคลื่นสี่เหลี่ยมจากฮาร์มอนิก}
\end{figure}

\subsubsection{7.6 ออร์โทกอนัลลิตีของฟังก์ชันไซน์และโคไซน์}

\paragraph{7.6.1 แนวคิดผลคูณภายใน}

เหตุผลที่สามารถแยกสัมประสิทธิ์ของฮาร์มอนิกแต่ละลำดับได้ คือ ฟังก์ชันไซน์และโคไซน์ที่มีความถี่ฮาร์มอนิกต่างกันมีสมบัติ \textbf{ออร์โทกอนัล (orthogonal)} บนช่วงหนึ่งคาบ เปรียบได้กับเวกเตอร์ตั้งฉากกันในเรขาคณิต

กำหนดผลคูณภายในของฟังก์ชันสองฟังก์ชันบนช่วงหนึ่งคาบเป็น

\[
\langle f,g\rangle=\int_{t_0}^{t_0+T}f(t)g(t)\,dt
\]

ฟังก์ชัน \(f\) และ \(g\) ออร์โทกอนัลกันเมื่อ

\[
\langle f,g\rangle=0
\]

\paragraph{7.6.2 ความสัมพันธ์ออร์โทกอนัลที่สำคัญ}

สำหรับจำนวนเต็มบวก \(m\) และ \(n\)

\[
\int_{t_0}^{t_0+T}\cos(m\omega_0t)\cos(n\omega_0t)\,dt=
\begin{cases}
0, & m\ne n\\
\dfrac{T}{2}, & m=n
\end{cases}
\]

\[
\int_{t_0}^{t_0+T}\sin(m\omega_0t)\sin(n\omega_0t)\,dt=
\begin{cases}
0, & m\ne n\\
\dfrac{T}{2}, & m=n
\end{cases}
\]

และ

\[
\int_{t_0}^{t_0+T}\sin(m\omega_0t)\cos(n\omega_0t)\,dt=0
\]

นอกจากนี้

\[
\int_{t_0}^{t_0+T}\sin(n\omega_0t)\,dt=0
\]

\[
\int_{t_0}^{t_0+T}\cos(n\omega_0t)\,dt=0
\]

\textbf{ที่มา:} ใช้เอกลักษณ์ผลคูณเป็นผลบวก เช่น

\[
\cos A\cos B=\frac{1}{2}\left[\cos(A-B)+\cos(A+B)\right]
\]

เมื่ออินทิเกรตตลอดจำนวนรอบเต็ม พื้นที่บวกและลบของไซน์หรือโคไซน์หักล้างกัน ยกเว้นกรณีความถี่เดียวกันซึ่งพจน์ \(\cos(A-B)\) กลายเป็น 1

\paragraph{7.6.3 ความหมายต่อการหาสัมประสิทธิ์}

หากคูณสมการอนุกรมฟูเรียร์ด้วย \(\cos(m\omega_0t)\) แล้วอินทิเกรตตลอดหนึ่งคาบ พจน์ที่มีความถี่ต่างจาก \(m\omega_0\) จะหายไปทั้งหมด เหลือเฉพาะพจน์ \(a_m\) จึงสามารถแยกสัมประสิทธิ์แต่ละตัวได้ หลักการเดียวกันใช้กับ \(b_m\)

% [รูปที่ 6: ออร์โทกอนัลลิตีของไซน์และโคไซน์บนหนึ่งคาบ]
% ประเภทภาพ: Graph / Conceptual Illustration
% ตำแหน่งที่แนะนำ: หลังหัวข้อ 7.6.3
% วัตถุประสงค์การเรียนรู้: ทำให้เห็นเชิงกราฟว่าอินทิกรัลของผลคูณฟังก์ชันต่างความถี่หักล้างเป็นศูนย์
% คำอธิบายภาพ: แสดงกราฟ $\sin(\omega_0t)\sin(2\omega_0t)$ และแรเงาพื้นที่บวกกับลบในหนึ่งคาบให้มีขนาดรวมเท่ากัน พร้อมกราฟ $\sin^2(\omega_0t)$ ที่มีพื้นที่สุทธิ $T/2$
% องค์ประกอบที่ต้องมี: ช่วง $0$ ถึง $T$ เส้นศูนย์ พื้นที่บวกและลบ สมการอินทิกรัล
% ข้อความหรือป้ายกำกับ: “different harmonics → integral = 0”, “same harmonic → integral = T/2”
% ภาษาของป้ายกำกับ: ไทยและอังกฤษ
% สัดส่วนภาพ: 16:9
% Image Generation Prompt: Two-panel educational graph explaining orthogonality. Panel A plots sin(ω0t)sin(2ω0t) over one period with positive and negative shaded areas canceling, labeled integral equals zero. Panel B plots sin²(ω0t) over one period with total area T/2. Clear axes, Thai-English labels, white background, clean engineering textbook style, 16:9.
% Graph Generation Prompt: Over $t\in[0,T]$ with $T=2\pi$, plot $\sin(t)\sin(2t)$ and $\sin^2(t)$. Shade positive and negative areas separately and annotate the integrals 0 and $\pi=T/2$.
% Alt Text: กราฟแสดงพื้นที่บวกและลบของผลคูณฮาร์มอนิกต่างลำดับหักล้างกัน ขณะที่กำลังสองของฮาร์มอนิกมีพื้นที่สุทธิไม่เป็นศูนย์

\begin{figure}[htbp]
    \centering
    \includegraphics[width=0.8\textwidth]{images/figure-06.png}
    \caption{ออร์โทกอนัลลิตีของไซน์และโคไซน์บนหนึ่งคาบ}
\end{figure}

\subsubsection{7.7 สัมประสิทธิ์ฟูเรียร์}

\paragraph{7.7.1 สัมประสิทธิ์สำหรับคาบทั่วไป \(T\)}

ให้ \(f(t)\) มีคาบ \(T\) และ \(\omega_0=2\pi/T\) สัมประสิทธิ์ฟูเรียร์คือ

\[
a_0=\frac{2}{T}\int_{t_0}^{t_0+T}f(t)\,dt
\]

\[
a_n=\frac{2}{T}\int_{t_0}^{t_0+T}f(t)\cos(n\omega_0t)\,dt,\qquad n\ge1
\]

\[
b_n=\frac{2}{T}\int_{t_0}^{t_0+T}f(t)\sin(n\omega_0t)\,dt,\qquad n\ge1
\]

ช่วงอินทิเกรตสามารถเลือกช่วงใดก็ได้ที่ยาวเท่ากับหนึ่งคาบ เช่น \([0,T]\), \([-T/2,T/2]\) หรือ \([t_0,t_0+T]\) แต่ต้องนิยามฟังก์ชันถูกต้องตลอดช่วงนั้น

\paragraph{7.7.2 ความหมายของ \(a_0\)}

ค่าเฉลี่ยของฟังก์ชันหนึ่งคาบคือ

\[
f_{\text{avg}}=\frac{1}{T}\int_{t_0}^{t_0+T}f(t)\,dt
\]

เปรียบเทียบกับนิยาม \(a_0\) จะได้

\[
f_{\text{avg}}=\frac{a_0}{2}
\]

ดังนั้นต้องระวังว่า \(a_0\) มิใช่ค่าเฉลี่ยโดยตรง แต่ค่าเฉลี่ยคือ \(a_0/2\) ภายใต้รูปแบบอนุกรมที่ใช้ในบทนี้

\paragraph{7.7.3 ที่มาของสูตร \(a_n\)}

เริ่มจาก

\[
f(t)=\frac{a_0}{2}+\sum_{n=1}^{\infty}\left[a_n\cos(n\omega_0t)+b_n\sin(n\omega_0t)\right]
\]

คูณทั้งสองข้างด้วย \(\cos(m\omega_0t)\) และอินทิเกรตตลอดหนึ่งคาบ

\[
\int f(t)\cos(m\omega_0t)\,dt
=\frac{a_0}{2}\int\cos(m\omega_0t)\,dt
+\sum a_n\int\cos(n\omega_0t)\cos(m\omega_0t)\,dt
+\sum b_n\int\sin(n\omega_0t)\cos(m\omega_0t)\,dt
\]

จากออร์โทกอนัลลิตี พจน์ทั้งหมดเป็นศูนย์ยกเว้นกรณี \(n=m\)

\[
\int_{t_0}^{t_0+T} f(t)\cos(m\omega_0t)\,dt=a_m\frac{T}{2}
\]

จึงได้

\[
a_m=\frac{2}{T}\int_{t_0}^{t_0+T} f(t)\cos(m\omega_0t)\,dt
\]

วิธีเดียวกันให้สูตร \(b_m\)

\paragraph{7.7.4 สมมติฐานและข้อจำกัด}

สูตรสัมประสิทธิ์ใช้ได้เมื่อฟังก์ชันสามารถอินทิเกรตได้บนหนึ่งคาบ ในงานวิศวกรรมมักใช้กับฟังก์ชันต่อเนื่องเป็นช่วงหรือมีจุดกระโดดจำนวนจำกัด หากรูปคลื่นมีอิมพัลส์หรือพฤติกรรมเชิงการแจกแจง ต้องใช้ทฤษฎีที่สูงขึ้นซึ่งอยู่นอกขอบเขตบทนี้

\paragraph{ตัวอย่างที่ 4: อนุกรมฟูเรียร์ของแรงดันคงที่}

\textbf{ตัวอย่างเพิ่มเติมเพื่อการเรียนรู้}\\
ให้ \(v(t)=V_0\) เป็นฟังก์ชันคงที่และกำหนดให้มีคาบ \(T\) จงหาสัมประสิทธิ์ฟูเรียร์

\textbf{1. วิเคราะห์สมมาตรและค่าเฉลี่ย}\\
สัญญาณมีค่าเท่ากันตลอดเวลา จึงควรมีเพียง DC

\textbf{2. คำนวณ \(a_0\)}

\[
a_0=\frac{2}{T}\int_{0}^{T}V_0\,dt
=\frac{2}{T}(V_0T)=2V_0
\]

ดังนั้น

\[
\frac{a_0}{2}=V_0
\]

\textbf{3. คำนวณ \(a_n\)}

\[
a_n=\frac{2V_0}{T}\int_0^T\cos(n\omega_0t)\,dt=0
\]

\textbf{4. คำนวณ \(b_n\)}

\[
b_n=\frac{2V_0}{T}\int_0^T\sin(n\omega_0t)\,dt=0
\]

\textbf{5. สรุป}

\[
v(t)=\frac{a_0}{2}=V_0
\]

\textbf{การตรวจสอบ:} ไม่มีส่วนที่เปลี่ยนตามเวลา จึงไม่มีองค์ประกอบ AC หรือฮาร์มอนิก

\subsubsection{7.8 อนุกรมฟูเรียร์ของฟังก์ชันคาบ}

\paragraph{7.8.1 รูปแบบคาบ \(2L\)}

ในตำราคณิตศาสตร์มักกำหนดฟังก์ชันบนช่วง \([-L,L]\) และให้มีคาบ \(2L\) ดังนั้น

\[
T=2L,\qquad \omega_0=\frac{\pi}{L}
\]

อนุกรมฟูเรียร์เป็น

\[
f(t)=\frac{a_0}{2}+\sum_{n=1}^{\infty}\left[a_n\cos\left(\frac{n\pi t}{L}\right)+b_n\sin\left(\frac{n\pi t}{L}\right)\right]
\]

สัมประสิทธิ์คือ

\[
a_0=\frac{1}{L}\int_{-L}^{L}f(t)\,dt
\]

\[
a_n=\frac{1}{L}\int_{-L}^{L}f(t)\cos\left(\frac{n\pi t}{L}\right)\,dt
\]

\[
b_n=\frac{1}{L}\int_{-L}^{L}f(t)\sin\left(\frac{n\pi t}{L}\right)\,dt
\]

รูปแบบนี้สมมูลกับสูตรคาบ \(T\) เพียงเปลี่ยนตัวแปร \(T=2L\)

\paragraph{7.8.2 ฟังก์ชันนิยามเป็นช่วง}

รูปคลื่นวิศวกรรมจำนวนมากนิยามเป็นช่วง เช่น

\[
f(t)=
\begin{cases}
V_m, & 0<t<T/2\\
-V_m, & -T/2<t<0
\end{cases}
\]

และทำซ้ำด้วยคาบ \(T\) การหาสัมประสิทธิ์ต้องแบ่งอินทิกรัลตามช่วงที่นิยามต่างกัน หรือใช้สมมาตรเพื่อลดจำนวนช่วง

\paragraph{7.8.3 ขั้นตอนทั่วไปในการหาสัมประสิทธิ์}

\begin{enumerate}
\def\labelenumi{\arabic{enumi}.}
\item
  ระบุคาบ \(T\) และหา \(\omega_0\)
\item
  เลือกช่วงอินทิเกรตหนึ่งคาบที่ทำให้ฟังก์ชันเขียนง่าย
\item
  ตรวจสอบค่าเฉลี่ยและสมมาตร
\item
  เขียน \(a_0\), \(a_n\), \(b_n\) ก่อนแทนฟังก์ชัน
\item
  แบ่งอินทิกรัลตามช่วงที่ฟังก์ชันเปลี่ยนนิยาม
\item
  คำนวณและจัดรูปให้เห็นพฤติกรรมตาม \(n\)
\item
  ตรวจว่าลำดับใดเป็นศูนย์ เช่น ฮาร์มอนิกคู่หรือฮาร์มอนิกคี่
\item
  เขียนอนุกรมและตรวจสอบหน่วย
\item
  ทดลองแทน \(n=1,2,3\) เพื่อดูความสมเหตุสมผล
\end{enumerate}

% [รูปที่ 7: กระบวนการหาสัมประสิทธิ์ฟูเรียร์อย่างเป็นระบบ]
% ประเภทภาพ: Flowchart
% ตำแหน่งที่แนะนำ: หลังหัวข้อ 7.8.3
% วัตถุประสงค์การเรียนรู้: ให้ผู้เรียนมีขั้นตอนแก้โจทย์ที่ตรวจสอบได้และลดข้อผิดพลาดจากการเลือกช่วงหรือใช้สูตรผิด
% คำอธิบายภาพ: ผังงานเริ่มจาก Define one period → Find $T,\omega_0$ → Check symmetry → Compute DC → Compute $a_n,b_n$ → Simplify by harmonic order → Write series → Verify waveform and units
% องค์ประกอบที่ต้องมี: จุดตัดสินใจ “Even?”, “Odd?”, ลูกศรไปยัง $b_n=0$ หรือ $a_0=a_n=0$ และกล่องตรวจสอบผล
% ข้อความหรือป้ายกำกับ: “กำหนดหนึ่งคาบ”, “ตรวจสมมาตร”, “หาค่าเฉลี่ย”, “อินทิเกรต”, “ตรวจฮาร์มอนิก”, “สังเคราะห์กลับ”
% ภาษาของป้ายกำกับ: ไทยและอังกฤษ
% สัดส่วนภาพ: 16:9
% Image Generation Prompt: Clean engineering flowchart for calculating Fourier series coefficients. Steps: define one period; determine T and ω0; choose integration interval; decision diamonds for even or odd symmetry; compute DC a0/2; compute an and bn; simplify odd/even harmonic pattern; write the series; verify units and reconstruct waveform. Thai-English labels, white background, clear arrows, textbook style, 16:9.
% Alt Text: ผังงานแสดงขั้นตอนตั้งแต่กำหนดคาบ ตรวจสมมาตร หาสัมประสิทธิ์ ไปจนถึงตรวจสอบการสังเคราะห์รูปคลื่น

\begin{figure}[htbp]
    \centering
    \includegraphics[width=0.8\textwidth]{images/figure-07.png}
    \caption{กระบวนการหาสัมประสิทธิ์ฟูเรียร์อย่างเป็นระบบ}
\end{figure}

\subsubsection{7.9 ฟังก์ชันคู่ ฟังก์ชันคี่ และการลดรูปอนุกรมฟูเรียร์}

\paragraph{7.9.1 ฟังก์ชันคู่}

ฟังก์ชัน \(f(t)\) เป็นฟังก์ชันคู่เมื่อ

\[
f(-t)=f(t)
\]

กราฟสมมาตรกับแกนตั้ง ตัวอย่าง ได้แก่ \(t^2\), \(|t|\), \(\cos t\)

สมบัติของผลคูณคือ

\begin{itemize}
\item
  คู่ \(\times\) คู่ = คู่
\item
  คี่ \(\times\) คี่ = คู่
\item
  คู่ \(\times\) คี่ = คี่
\end{itemize}

อินทิกรัลของฟังก์ชันคี่บนช่วงสมมาตรเป็นศูนย์

\[
\int_{-L}^{L}g_{\text{odd}}(t)\,dt=0
\]

ถ้า \(f(t)\) เป็นฟังก์ชันคู่ เนื่องจาก \(\sin(n\pi t/L)\) เป็นฟังก์ชันคี่ ผลคูณ \(f(t)\sin(n\pi t/L)\) เป็นคี่ จึงได้

\[
b_n=0
\]

อนุกรมเหลือเฉพาะ DC และโคไซน์

\[
f(t)=\frac{a_0}{2}+\sum_{n=1}^{\infty}a_n\cos\left(\frac{n\pi t}{L}\right)
\]

และสามารถลดช่วงอินทิเกรตได้

\[
a_0=\frac{2}{L}\int_0^L f(t)\,dt
\]

\[
a_n=\frac{2}{L}\int_0^L f(t)\cos\left(\frac{n\pi t}{L}\right)\,dt
\]

\paragraph{7.9.2 ฟังก์ชันคี่}

ฟังก์ชัน \(f(t)\) เป็นฟังก์ชันคี่เมื่อ

\[
f(-t)=-f(t)
\]

กราฟมีสมมาตรแบบหมุน 180 องศารอบจุดกำเนิด ตัวอย่าง ได้แก่ \(t\), \(t^3\), \(\sin t\)

ถ้า \(f(t)\) เป็นคี่ ค่าเฉลี่ยบนช่วงสมมาตรเป็นศูนย์ และผลคูณกับโคไซน์เป็นฟังก์ชันคี่ จึงได้

\[
a_0=0,\qquad a_n=0
\]

อนุกรมเหลือเฉพาะไซน์

\[
f(t)=\sum_{n=1}^{\infty}b_n\sin\left(\frac{n\pi t}{L}\right)
\]

โดย

\[
b_n=\frac{2}{L}\int_0^L f(t)\sin\left(\frac{n\pi t}{L}\right)\,dt
\]

\paragraph{7.9.3 ฟังก์ชันที่ไม่เป็นทั้งคู่และคี่}

ฟังก์ชันทั่วไปอาจแยกเป็นส่วนคู่และส่วนคี่ได้

\[
f_e(t)=\frac{f(t)+f(-t)}{2}
\]

\[
f_o(t)=\frac{f(t)-f(-t)}{2}
\]

และ

\[
f(t)=f_e(t)+f_o(t)
\]

ส่วนคู่สร้างสัมประสิทธิ์โคไซน์ ส่วนคี่สร้างสัมประสิทธิ์ไซน์ แนวคิดนี้ช่วยอธิบายว่าทำไมสัญญาณทั่วไปจึงต้องมีทั้ง \(a_n\) และ \(b_n\)

\paragraph{7.9.4 สมมาตรครึ่งคลื่น}

หากฟังก์ชันมีสมมาตรครึ่งคลื่น

\[
f\left(t+\frac{T}{2}\right)=-f(t)
\]

องค์ประกอบ DC เป็นศูนย์และฮาร์มอนิกคู่หายไป เหลือเฉพาะฮาร์มอนิกคี่ อย่างไรก็ตาม สมมาตรครึ่งคลื่นไม่เท่ากับฟังก์ชันคี่เสมอไป เพราะขึ้นกับตำแหน่งจุดกำเนิดเวลา

\paragraph{7.9.5 สมมาตรหนึ่งในสี่คลื่น}

รูปคลื่นบางชนิดมีทั้งสมมาตรคู่หรือคี่และสมมาตรครึ่งคลื่น ทำให้สามารถคำนวณจากช่วงหนึ่งในสี่คาบได้ การใช้สมมาตรนี้ช่วยลดการอินทิเกรต แต่ต้องตรวจสอบรูปคลื่นและจุดกำเนิดอย่างระมัดระวัง

% [รูปที่ 8: การจำแนกสมมาตรของฟังก์ชันและผลต่อสัมประสิทธิ์]
% ประเภทภาพ: Infographic / Graph
% ตำแหน่งที่แนะนำ: หลังหัวข้อ 7.9.5
% วัตถุประสงค์การเรียนรู้: ช่วยให้เลือกได้ทันทีว่าสัมประสิทธิ์ชุดใดเป็นศูนย์
% คำอธิบายภาพ: ตารางภาพ 4 ช่อง ได้แก่ even, odd, half-wave symmetry และ no symmetry แต่ละช่องมีกราฟตัวอย่างและผลสรุปต่อ $a_0,a_n,b_n$ และ harmonic order
% องค์ประกอบที่ต้องมี: เส้นแกน $t=0$, สัญลักษณ์สะท้อน/หมุน, กล่องสูตร $b_n=0$, $a_0=a_n=0$, even harmonics = 0
% ข้อความหรือป้ายกำกับ: “Even”, “Odd”, “Half-wave symmetry”, “General waveform”
% ภาษาของป้ายกำกับ: ไทยและอังกฤษ
% สัดส่วนภาพ: 16:9
% Image Generation Prompt: Four-panel engineering mathematics infographic classifying waveform symmetry. Panel 1 even symmetry with mirror about vertical axis and result bn=0. Panel 2 odd symmetry with rotational symmetry and result a0=an=0. Panel 3 half-wave symmetry with second half inverted and result DC and even harmonics vanish. Panel 4 general waveform requiring both sine and cosine coefficients. Thai-English labels, clean vector graphs, white background, 16:9.
% Alt Text: อินโฟกราฟิกแสดงฟังก์ชันคู่ คี่ สมมาตรครึ่งคลื่น และฟังก์ชันทั่วไป พร้อมผลต่อสัมประสิทธิ์ฟูเรียร์

\begin{figure}[htbp]
    \centering
    \includegraphics[width=0.8\textwidth]{images/figure-08.png}
    \caption{การจำแนกสมมาตรของฟังก์ชันและผลต่อสัมประสิทธิ์}
\end{figure}

\paragraph{ตัวอย่างที่ 5: อนุกรมฟูเรียร์ของรูปคลื่นสี่เหลี่ยมสมมาตรคี่}

\textbf{ตัวอย่างเพิ่มเติมเพื่อการเรียนรู้}\\
ให้แรงดันคาบ \(T\) นิยามบนช่วง \(-T/2<t<T/2\) เป็น

\[
v(t)=
\begin{cases}
-V_m, & -T/2<t<0\\
V_m, & 0<t<T/2
\end{cases}
\]

และทำซ้ำทุกคาบ จงหาอนุกรมฟูเรียร์

\textbf{1. วิเคราะห์รูปคลื่น}\\
มี \(v(-t)=-v(t)\) จึงเป็นฟังก์ชันคี่

ดังนั้น

\[
a_0=0,\qquad a_n=0
\]

เหลือเพียง \(b_n\)

\textbf{2. เขียนสูตรโดยใช้สมมาตร}

\[
b_n=\frac{4}{T}\int_0^{T/2}V_m\sin(n\omega_0t)\,dt
\]

เนื่องจาก \(\omega_0=2\pi/T\)

\textbf{3. อินทิเกรต}

\[
b_n=\frac{4V_m}{T}\left[-\frac{\cos(n\omega_0t)}{n\omega_0}\right]_0^{T/2}
\]

\[
b_n=\frac{4V_m}{Tn\omega_0}\left[1-\cos\left(n\omega_0\frac{T}{2}\right)\right]
\]

แต่

\[
n\omega_0\frac{T}{2}=n\pi
\]

และ \(T\omega_0=2\pi\) จึงได้

\[
b_n=\frac{2V_m}{n\pi}\left[1-(-1)^n\right]
\]

\textbf{4. พิจารณาลำดับ \(n\)}

\begin{itemize}
\item
  ถ้า \(n\) เป็นจำนวนคู่ \((-1)^n=1\) จึง \(b_n=0\)
\item
  ถ้า \(n\) เป็นจำนวนคี่ \((-1)^n=-1\) จึง
\end{itemize}

\[
b_n=\frac{4V_m}{n\pi}
\]

\textbf{5. เขียนอนุกรม}

\[
v(t)=\frac{4V_m}{\pi}\left[\sin(\omega_0t)+\frac{1}{3}\sin(3\omega_0t)+\frac{1}{5}\sin(5\omega_0t)+\cdots\right]
\]

\textbf{6. ตรวจสอบ}

\begin{itemize}
\item
  ไม่มี DC เพราะพื้นที่บวกและลบเท่ากัน
\item
  ไม่มีโคไซน์เพราะเป็นฟังก์ชันคี่
\item
  ไม่มีฮาร์มอนิกคู่เพราะมีสมมาตรครึ่งคลื่น
\item
  ขนาดฮาร์มอนิกลดลงตาม \(1/n\)
\end{itemize}

\textbf{ความหมายทางกายภาพ:} ขอบที่เปลี่ยนอย่างฉับพลันต้องใช้องค์ประกอบความถี่สูงจำนวนมากเพื่อสร้างรูปคลื่นให้คม

\paragraph{ตัวอย่างที่ 6: ขบวนพัลส์สมมาตรคู่}

\textbf{ตัวอย่างเพิ่มเติมเพื่อการเรียนรู้}\\
ให้สัญญาณพัลส์คาบ \(T\) มีแอมพลิจูด \(A\) และความกว้างพัลส์ \(\tau\) โดยกำหนดหนึ่งคาบบนช่วง \(-T/2<t<T/2\) เป็น

\[
x(t)=
\begin{cases}
A, & |t|<\tau/2\\
0, & \tau/2<|t|<T/2
\end{cases}
\]

จงหาสัมประสิทธิ์ฟูเรียร์

\textbf{1. วิเคราะห์สมมาตร}\\
\(x(-t)=x(t)\) จึงเป็นฟังก์ชันคู่ และ

\[
b_n=0
\]

\textbf{2. หาองค์ประกอบ DC}

\[
a_0=\frac{2}{T}\int_{-T/2}^{T/2}x(t)\,dt
\]

ฟังก์ชันไม่เป็นศูนย์เฉพาะช่วงยาว \(\tau\)

\[
a_0=\frac{2A\tau}{T}
\]

ดังนั้นค่าเฉลี่ยคือ

\[
\frac{a_0}{2}=A\frac{\tau}{T}
\]

กำหนดอัตราส่วนหน้าที่ \(D=\tau/T\) จะได้ค่าเฉลี่ย \(AD\)

\textbf{3. หา \(a_n\)}

\[
a_n=\frac{2}{T}\int_{-\tau/2}^{\tau/2}A\cos(n\omega_0t)\,dt
\]

เพราะอินทิกรัลเป็นฟังก์ชันคู่

\[
a_n=\frac{4A}{T}\int_0^{\tau/2}\cos(n\omega_0t)\,dt
\]

\[
a_n=\frac{4A}{Tn\omega_0}\sin\left(\frac{n\omega_0\tau}{2}\right)
\]

ใช้ \(\omega_0=2\pi/T\)

\[
a_n=\frac{2A}{n\pi}\sin\left(n\pi\frac{\tau}{T}\right)
\]

หรือ

\[
a_n=\frac{2A}{n\pi}\sin(n\pi D)
\]

\textbf{4. เขียนอนุกรม}

\[
x(t)=AD+\sum_{n=1}^{\infty}\frac{2A}{n\pi}\sin(n\pi D)\cos(n\omega_0t)
\]

\textbf{5. แปลความหมาย}

\begin{itemize}
\item
  ค่าเฉลี่ยขึ้นกับพื้นที่พัลส์ \(A\tau\) ต่อหนึ่งคาบ
\item
  การเปลี่ยน duty cycle ทำให้ทั้ง DC และรูปแบบสเปกตรัมเปลี่ยน
\item
  สัมประสิทธิ์บางลำดับเป็นศูนย์เมื่อ \(\sin(n\pi D)=0\)
\end{itemize}

\subsubsection{7.10 สเปกตรัมของสัญญาณและฮาร์มอนิก}

\paragraph{7.10.1 รูปขนาด--เฟส}

พจน์ไซน์และโคไซน์ของฮาร์มอนิกเดียวกันสามารถรวมเป็นโคไซน์ที่มีขนาดและเฟสได้

\[
a_n\cos(n\omega_0t)+b_n\sin(n\omega_0t)
=A_n\cos(n\omega_0t-\theta_n)
\]

โดย

\[
A_n=\sqrt{a_n^2+b_n^2}
\]

และ

\[
\theta_n=\operatorname{atan2}(b_n,a_n)
\]

ดังนั้น

\[
f(t)=\frac{a_0}{2}+\sum_{n=1}^{\infty}A_n\cos(n\omega_0t-\theta_n)
\]

\textbf{ข้อควรระวัง:} สูตรเฟสขึ้นกับรูปแบบที่เลือก หากเขียน \(A_n\cos(n\omega_0t+\phi_n)\) จะได้ \(\phi_n=-\theta_n\) จึงต้องระบุเครื่องหมายในสมการให้ชัดเจนและใช้ให้สม่ำเสมอ

\paragraph{7.10.2 สเปกตรัมขนาด}

สเปกตรัมขนาดแสดง \(A_n\) ที่ความถี่ \(nf_0\) เป็นแท่งไม่ต่อเนื่อง เพราะฟังก์ชันเป็นคาบมีองค์ประกอบเฉพาะความถี่ฮาร์มอนิก ส่วน DC อยู่ที่ 0 Hz และมีขนาด \(|a_0/2|\)

\paragraph{7.10.3 สเปกตรัมเฟส}

สเปกตรัมเฟสแสดง \(\theta_n\) ของแต่ละฮาร์มอนิก เฟสมีบทบาทต่อการจัดตำแหน่งและรูปร่างของสัญญาณในเวลา แม้สัญญาณสองสัญญาณมีสเปกตรัมขนาดเท่ากัน แต่ถ้าเฟสต่างกัน รูปคลื่นในโดเมนเวลาอาจต่างกันอย่างมาก

\paragraph{7.10.4 รูปเชิงซ้อน}

อนุกรมฟูเรียร์ยังเขียนด้วยเลขชี้กำลังเชิงซ้อนได้

\[
f(t)=\sum_{n=-\infty}^{\infty}C_ne^{jn\omega_0t}
\]

โดย

\[
C_n=\frac{1}{T}\int_{t_0}^{t_0+T}f(t)e^{-jn\omega_0t}\,dt
\]

ความสัมพันธ์กับสัมประสิทธิ์ตรีโกณมิติสำหรับ \(n>0\) คือ

\[
C_n=\frac{a_n-jb_n}{2}
\]

\[
C_{-n}=\frac{a_n+jb_n}{2}
\]

และ

\[
C_0=\frac{a_0}{2}
\]

สำหรับสัญญาณจริง

\[
C_{-n}=C_n^*
\]

เรียกว่า \textbf{สมมาตรคอนจูเกต (conjugate symmetry)} รูปเชิงซ้อนมีประโยชน์มากในวิชาสัญญาณและระบบ แต่บทนี้เน้นรูปตรีโกณมิติเป็นหลัก

% [รูปที่ 9: สเปกตรัมขนาดและเฟสของรูปคลื่นเป็นคาบ]
% ประเภทภาพ: Graph
% ตำแหน่งที่แนะนำ: หลังหัวข้อ 7.10.4
% วัตถุประสงค์การเรียนรู้: เชื่อมโยงสัมประสิทธิ์ $a_n,b_n$ กับแอมพลิจูดและเฟสที่แต่ละฮาร์มอนิก
% คำอธิบายภาพ: แผงบนเป็นรูปคลื่นในเวลา แผงกลางเป็นสเปกตรัมขนาดที่ 0, $f_0$, $2f_0$, ... แผงล่างเป็นสเปกตรัมเฟส แสดงเฉพาะจุดที่ขนาดไม่เป็นศูนย์
% องค์ประกอบที่ต้องมี: แกนความถี่ Hz แท่ง DC และฮาร์มอนิก ป้าย $A_n$ และ $\theta_n$
% ข้อความหรือป้ายกำกับ: “Time Domain”, “Magnitude Spectrum”, “Phase Spectrum”
% ภาษาของป้ายกำกับ: ไทยและอังกฤษ
% สัดส่วนภาพ: 16:9
% Image Generation Prompt: Three-tier engineering signal illustration. Top: a periodic nonsinusoidal waveform in time. Middle: discrete magnitude line spectrum with DC and harmonics at f0, 2f0, 3f0, etc. Bottom: corresponding phase spectrum at the same frequencies. Label A_n and theta_n, frequency in hertz, clean white-background textbook style, Thai-English labels, 16:9.
% Graph Generation Prompt: Use example coefficients $a_0/2=2$, $(a_1,b_1)=(3,4)$, $(a_2,b_2)=(0,2)$, $(a_3,b_3)=(-1,0)$ with $f_0=50$ Hz. Plot magnitude $A_n$ and phase $\theta_n=\operatorname{atan2}(b_n,a_n)$ at 0, 50, 100, and 150 Hz.
% Alt Text: รูปคลื่นในเวลาเชื่อมกับแท่งสเปกตรัมขนาดและเฟสที่ความถี่มูลฐานและฮาร์มอนิก

\begin{figure}[htbp]
    \centering
    \includegraphics[width=0.8\textwidth]{images/figure-09.png}
    \caption{สเปกตรัมขนาดและเฟสของรูปคลื่นเป็นคาบ}
\end{figure}

\paragraph{ตัวอย่างที่ 7: แปลงสัมประสิทธิ์เป็นขนาดและเฟส}

\textbf{ตัวอย่างเพิ่มเติมเพื่อการเรียนรู้}\\
ฮาร์มอนิกลำดับหนึ่งของสัญญาณมี \(a_1=3\ \text{V}\) และ \(b_1=4\ \text{V}\) จงเขียนในรูปขนาด--เฟส

\textbf{1. ขนาด}

\[
A_1=\sqrt{3^2+4^2}=5\ \text{V}
\]

\textbf{2. เฟส}

\[
\theta_1=\operatorname{atan2}(4,3)=53.13^\circ
\]

\textbf{3. เขียนพจน์}

\[
3\cos(\omega_0t)+4\sin(\omega_0t)
=5\cos(\omega_0t-53.13^\circ)
\]

\textbf{4. ตรวจสอบ}

\[
5\cos 53.13^\circ=3,
\qquad
5\sin 53.13^\circ=4
\]

จึงถูกต้อง

\subsubsection{7.11 ความหมายทางกายภาพขององค์ประกอบความถี่ในสัญญาณไฟฟ้า}

\paragraph{7.11.1 องค์ประกอบ DC}

องค์ประกอบ \(a_0/2\) คือค่าเฉลี่ยของรูปคลื่น หากแรงดันสลับมี DC offset ค่าเฉลี่ยจะไม่เป็นศูนย์ ในวงจรบางชนิด DC อาจทำให้แกนแม่เหล็กมีไบอัส ทำให้อุปกรณ์วัดอ่านค่าเปลี่ยน หรือทำให้สัญญาณอยู่นอกช่วงทำงานที่ต้องการ

\paragraph{7.11.2 ฮาร์มอนิกมูลฐาน}

ฮาร์มอนิกลำดับที่ 1 มีความถี่ \(f_0\) และมักเป็นองค์ประกอบหลักที่สัมพันธ์กับการทำงานพื้นฐานของระบบ เช่น ระบบไฟฟ้า 50 Hz หรือสัญญาณอ้างอิงของวงจรควบคุม

\paragraph{7.11.3 ฮาร์มอนิกลำดับสูง}

ฮาร์มอนิกลำดับสูงบ่งชี้รายละเอียดที่เปลี่ยนเร็วของรูปคลื่น เช่น ขอบชันหรือมุมแหลม ยิ่งรูปคลื่นเปลี่ยนฉับพลัน สัมประสิทธิ์มักลดลงช้ากว่าและต้องใช้ความถี่สูงมากขึ้นในการประมาณรูปคลื่น

ความถี่สูงมีผลแตกต่างจากความถี่ต่ำในวงจร เนื่องจากรีแอกแตนซ์ของตัวเหนี่ยวนำและตัวเก็บประจุขึ้นกับความถี่

\[
X_L=2\pi fL
\]

\[
X_C=\frac{1}{2\pi fC}
\]

ดังนั้นแต่ละฮาร์มอนิกอาจถูกวงจรลดทอนหรือขยายไม่เท่ากัน แม้บทนี้ยังไม่วิเคราะห์การตอบสนองของวงจรอย่างละเอียด แต่แนวคิดนี้อธิบายว่าทำไมการแยกสัญญาณตามความถี่จึงมีประโยชน์

\paragraph{7.11.4 เฟสและรูปร่างคลื่น}

เฟสกำหนดตำแหน่งสัมพัทธ์ของยอดและจุดตัดศูนย์ การเปลี่ยนเฟสฮาร์มอนิกโดยไม่เปลี่ยนขนาดสามารถเปลี่ยนรูปร่างคลื่นได้ จึงไม่ควรสรุปคุณลักษณะของสัญญาณจากสเปกตรัมขนาดเพียงอย่างเดียวในทุกกรณี

\paragraph{7.11.5 การซ้อนทับในระบบเชิงเส้น}

หากระบบเป็นเชิงเส้น สามารถพิจารณาการตอบสนองต่อแต่ละฮาร์มอนิกแยกกัน แล้วรวมผลตอบสนองได้ หลักการนี้เป็นหัวใจของการวิเคราะห์วงจรและระบบในโดเมนความถี่ แต่ถ้าระบบไม่เชิงเส้น ระบบอาจสร้างความถี่ใหม่ที่ไม่ปรากฏในอินพุต จึงต้องใช้การวิเคราะห์เพิ่มเติม

\subsubsection{7.12 การลู่เข้าของอนุกรมฟูเรียร์ เงื่อนไขดิริชเลต์ และปรากฏการณ์กิบส์}

\paragraph{7.12.1 เงื่อนไขดิริชเลต์ในเชิงใช้งาน}

เงื่อนไขที่นิยมใช้เป็นแนวทางเพียงพอสำหรับการลู่เข้าของอนุกรมฟูเรียร์ในงานพื้นฐาน ได้แก่ ในหนึ่งคาบ ฟังก์ชัน

\begin{enumerate}
\def\labelenumi{\arabic{enumi}.}
\item
  สามารถอินทิเกรตค่าสัมบูรณ์ได้
\item
  มีจำนวนจุดสูงสุดและต่ำสุดจำกัด
\item
  มีจำนวนจุดไม่ต่อเนื่องจำกัด และค่าซ้าย--ขวาจำกัด
\end{enumerate}

รูปคลื่นไฟฟ้าทั่วไปที่ต่อเนื่องเป็นช่วง เช่น สี่เหลี่ยม สามเหลี่ยม และฟันเลื่อย มักผ่านเงื่อนไขเหล่านี้

\paragraph{7.12.2 ค่าที่จุดต่อเนื่องและจุดกระโดด}

ที่จุดซึ่ง \(f(t)\) ต่อเนื่อง อนุกรมฟูเรียร์ลู่เข้าหา \(f(t)\) ภายใต้เงื่อนไขที่เหมาะสม

ที่จุดกระโดด \(t_d\) อนุกรมลู่เข้าหาค่าเฉลี่ยของค่าด้านซ้ายและด้านขวา

\[
S(t_d)=\frac{f(t_d^-)+f(t_d^+)}{2}
\]

ตัวอย่างเช่น รูปคลื่นสี่เหลี่ยมที่กระโดดจาก \(-V_m\) ไป \(V_m\) จะมีผลรวมฟูเรียร์ที่จุดกระโดดเท่ากับ 0 แม้จะกำหนดค่าฟังก์ชันที่จุดเดียวเป็นค่าอื่นก็ตาม

\paragraph{7.12.3 ปรากฏการณ์กิบส์}

เมื่อใช้ผลรวมแบบตัดทอนประมาณฟังก์ชันที่มีจุดกระโดด จะเกิดการสั่นและพุ่งเกินใกล้จุดไม่ต่อเนื่อง เรียกว่า \textbf{ปรากฏการณ์กิบส์} เมื่อเพิ่มจำนวนพจน์ ช่วงที่เกิดการสั่นจะแคบลง แต่สัดส่วนการพุ่งเกินสูงสุดไม่หายไปทั้งหมด โดยมีค่าประมาณร้อยละ 9 ของขนาดการกระโดด

ประเด็นสำคัญคือ การเพิ่มฮาร์มอนิกมิได้ทำให้การประมาณ ``ผิดมากขึ้น'' พลังงานของความคลาดเคลื่อนโดยรวมลดลงและบริเวณผิดพลาดแคบลง แต่ค่าพุ่งเกินเฉพาะบริเวณจุดกระโดดยังคงปรากฏ

\paragraph{7.12.4 ความราบเรียบกับอัตราการลดลงของสัมประสิทธิ์}

โดยทั่วไป รูปคลื่นที่ราบเรียบกว่าจะมีสัมประสิทธิ์ฮาร์มอนิกลดลงเร็วกว่า

\begin{itemize}
\item
  รูปคลื่นมีจุดกระโดด เช่น สี่เหลี่ยม: สัมประสิทธิ์มักลดระดับประมาณ \(1/n\)
\item
  รูปคลื่นต่อเนื่องแต่มีมุม เช่น สามเหลี่ยม: สัมประสิทธิ์มักลดระดับประมาณ \(1/n^2\)
\item
  ฟังก์ชันที่ราบเรียบมาก: สัมประสิทธิ์อาจลดลงรวดเร็วกว่านั้น
\end{itemize}

ความสัมพันธ์นี้มีความหมายเชิงวิศวกรรมว่า สัญญาณที่มีขอบชันต้องการแบนด์วิดท์สูงกว่าในการส่งผ่านหรือสร้างซ้ำให้คงรูปร่าง

% [รูปที่ 10: ปรากฏการณ์กิบส์บริเวณจุดไม่ต่อเนื่อง]
% ประเภทภาพ: Graph
% ตำแหน่งที่แนะนำ: หลังหัวข้อ 7.12.4
% วัตถุประสงค์การเรียนรู้: อธิบายว่าการเพิ่มจำนวนพจน์ทำให้บริเวณสั่นแคบลงแต่ค่าพุ่งเกินไม่หายทั้งหมด
% คำอธิบายภาพ: แสดงรูปคลื่นสี่เหลี่ยมอุดมคติและผลรวมฟูเรียร์ $N=5,15,50$ ซูมบริเวณจุดกระโดด ทำเครื่องหมาย overshoot ประมาณ 9% ของขนาดการกระโดด
% องค์ประกอบที่ต้องมี: จุดกระโดด เส้นระดับก่อนและหลัง เส้นผลรวมหลายค่า $N$ ลูกศร overshoot และคำอธิบาย “narrower but persistent”
% ข้อความหรือป้ายกำกับ: “Gibbs phenomenon”, “jump size”, “approximately 9% overshoot”
% ภาษาของป้ายกำกับ: ไทยและอังกฤษ
% สัดส่วนภาพ: 16:9
% Image Generation Prompt: Engineering mathematics graph illustrating Gibbs phenomenon near a square-wave discontinuity. Overlay ideal step and Fourier partial sums for N=5, 15, and 50. Include a zoomed inset showing overshoot about 9 percent of the jump, with the oscillation region becoming narrower as N increases but peak overshoot persisting. Clean white background, precise axes, Thai-English labels, 16:9.
% Graph Generation Prompt: Plot the odd square-wave Fourier partial sums with 5, 15, and 50 odd harmonics over a small interval around $t=0$, with ideal levels -1 and +1. Mark the maximum overshoot near 1.18 and explain that the jump size is 2, so the overshoot above +1 is about 0.18.
% Alt Text: ผลรวมฟูเรียร์ของรูปคลื่นสี่เหลี่ยมเกิดการพุ่งเกินใกล้จุดกระโดด โดยพื้นที่สั่นแคบลงเมื่อเพิ่มจำนวนฮาร์มอนิก

\begin{figure}[htbp]
    \centering
    \includegraphics[width=0.8\textwidth]{images/figure-10.png}
    \caption{ปรากฏการณ์กิบส์บริเวณจุดไม่ต่อเนื่อง}
\end{figure}

\paragraph{7.12.5 ความสัมพันธ์ของพาร์เซวาล}

ความสัมพันธ์ของพาร์เซวาลเชื่อมพลังงานกำลังสองเฉลี่ยในโดเมนเวลากับผลรวมกำลังสองของสัมประสิทธิ์ สำหรับฟังก์ชันจริงคาบ \(T\)

\[
\frac{1}{T}\int_{t_0}^{t_0+T}|f(t)|^2\,dt
=\frac{a_0^2}{4}+\frac{1}{2}\sum_{n=1}^{\infty}\left(a_n^2+b_n^2\right)
\]

หาก \(f(t)\) เป็นแรงดัน ด้านซ้ายเท่ากับ \(V_{\mathrm{rms}}^2\) ดังนั้นกำลังสองเฉลี่ยของรูปคลื่นเท่ากับผลรวมส่วน DC และส่วนของฮาร์มอนิกแต่ละลำดับ

\textbf{เงื่อนไขการใช้:} ฟังก์ชันต้องมีค่ากำลังสองที่อินทิเกรตได้บนหนึ่งคาบ

\textbf{ความหมายทางกายภาพ:} ฮาร์มอนิกแต่ละลำดับมีส่วนต่อค่ากำลังสองเฉลี่ยของสัญญาณอย่างเป็นอิสระ เพราะฟังก์ชันฐานออร์โทกอนัลกัน

\textbf{ขอบเขตของบท:} ใช้ความสัมพันธ์นี้เพื่อสร้างความเชื่อมโยงเบื้องต้นเท่านั้น การคำนวณ RMS กำลังไฟฟ้า และ THD จะศึกษาอย่างละเอียดในบทที่ 5

\subsubsection{7.13 ตัวอย่างโจทย์และการคำนวณแบบเป็นขั้นตอน}

หัวข้อนี้รวบรวมตัวอย่างที่เชื่อมแนวคิดหลายส่วนเข้าด้วยกัน โดยใช้ขั้นตอนวิเคราะห์ กำหนดตัวแปร เลือกสูตร คำนวณ ตรวจสอบ และแปลความหมาย

\paragraph{ตัวอย่างที่ 8: อนุกรมฟูเรียร์ของรูปคลื่นฟันเลื่อยสมมาตรคี่}

\textbf{ตัวอย่างเพิ่มเติมเพื่อการเรียนรู้}\\
ให้ฟังก์ชันคาบ \(2L\) นิยามบนช่วง \(-L<t<L\) เป็น

\[
f(t)=\frac{V_m}{L}t
\]

จงหาอนุกรมฟูเรียร์

\textbf{1. วิเคราะห์ฟังก์ชัน}

\[
f(-t)=\frac{V_m}{L}(-t)=-f(t)
\]

จึงเป็นฟังก์ชันคี่

\[
a_0=0,\qquad a_n=0
\]

\textbf{2. เลือกสูตร}

สำหรับคาบ \(2L\)

\[
b_n=\frac{2}{L}\int_0^L f(t)\sin\left(\frac{n\pi t}{L}\right)\,dt
\]

แทน \(f(t)=V_mt/L\)

\[
b_n=\frac{2V_m}{L^2}\int_0^L t\sin\left(\frac{n\pi t}{L}\right)\,dt
\]

\textbf{3. อินทิเกรตโดยส่วน}

กำหนด

\[
u=t,\qquad dv=\sin\left(\frac{n\pi t}{L}\right)dt
\]

ดังนั้น

\[
du=dt,
\qquad
v=-\frac{L}{n\pi}\cos\left(\frac{n\pi t}{L}\right)
\]

จึงได้

\[
\int_0^L t\sin\left(\frac{n\pi t}{L}\right)dt
=
\left[-\frac{Lt}{n\pi}\cos\left(\frac{n\pi t}{L}\right)\right]_0^L
+\frac{L}{n\pi}\int_0^L\cos\left(\frac{n\pi t}{L}\right)dt
\]

อินทิกรัลพจน์หลังเป็นศูนย์ และพจน์ขอบเขตให้

\[
\int_0^L t\sin\left(\frac{n\pi t}{L}\right)dt
=-\frac{L^2}{n\pi}(-1)^n
\]

\textbf{4. หาสัมประสิทธิ์}

\[
b_n=\frac{2V_m}{L^2}\left[-\frac{L^2}{n\pi}(-1)^n\right]
\]

\[
b_n=\frac{2V_m}{n\pi}(-1)^{n+1}
\]

\textbf{5. เขียนอนุกรม}

\[
f(t)=\frac{2V_m}{\pi}\sum_{n=1}^{\infty}\frac{(-1)^{n+1}}{n}
\sin\left(\frac{n\pi t}{L}\right)
\]

หรือเขียนพจน์แรก ๆ เป็น

\[
f(t)=\frac{2V_m}{\pi}\left[
\sin\left(\frac{\pi t}{L}\right)
-\frac{1}{2}\sin\left(\frac{2\pi t}{L}\right)
+\frac{1}{3}\sin\left(\frac{3\pi t}{L}\right)
-\cdots
\right]
\]

\textbf{6. ตรวจสอบ}

\begin{itemize}
\item
  ฟังก์ชันคี่จึงมีเฉพาะพจน์ไซน์
\item
  มีทั้งฮาร์มอนิกคู่และคี่ เพราะรูปคลื่นนี้ไม่มีสมมาตรครึ่งคลื่นตามนิยาม \(f(t+L)=-f(t)\) เมื่อพิจารณาจุดกระโดดของการต่อคาบอย่างระมัดระวัง รูปแบบสัมประสิทธิ์สลับเครื่องหมายและไม่ตัดฮาร์มอนิกคู่ทั้งหมด
\item
  สัมประสิทธิ์ลดตาม \(1/n\) เพราะการต่อคาบมีจุดกระโดดจาก \(V_m\) กลับไป \(-V_m\)
\end{itemize}

\paragraph{ตัวอย่างที่ 9: การประมาณรูปคลื่นสี่เหลี่ยมด้วยสามฮาร์มอนิกแรก}

\textbf{ตัวอย่างเพิ่มเติมเพื่อการเรียนรู้}\\
รูปคลื่นสี่เหลี่ยมสมมาตรคี่มี

\[
v(t)=\frac{4V_m}{\pi}\left[\sin(\omega_0t)+\frac{1}{3}\sin(3\omega_0t)+\frac{1}{5}\sin(5\omega_0t)+\cdots\right]
\]

จงประมาณค่า \(v(T/8)\) โดยใช้ฮาร์มอนิกคี่สามลำดับแรก

\textbf{1. กำหนดมุม}

\[
\omega_0\frac{T}{8}=\frac{2\pi}{T}\frac{T}{8}=\frac{\pi}{4}
\]

\textbf{2. เขียนผลรวมตัดทอน}

\[
S_3\left(\frac{T}{8}\right)
=\frac{4V_m}{\pi}\left[
\sin\left(\frac{\pi}{4}\right)
+\frac{1}{3}\sin\left(\frac{3\pi}{4}\right)
+\frac{1}{5}\sin\left(\frac{5\pi}{4}\right)
\right]
\]

\textbf{3. แทนค่าตรีโกณมิติ}

\[
\sin\left(\frac{\pi}{4}\right)=\frac{\sqrt{2}}{2}
\]

\[
\sin\left(\frac{3\pi}{4}\right)=\frac{\sqrt{2}}{2}
\]

\[
\sin\left(\frac{5\pi}{4}\right)=-\frac{\sqrt{2}}{2}
\]

ดังนั้น

\[
S_3\left(\frac{T}{8}\right)
=\frac{4V_m}{\pi}\frac{\sqrt{2}}{2}
\left(1+\frac{1}{3}-\frac{1}{5}\right)
\]

\[
S_3\left(\frac{T}{8}\right)
\approx1.020\,V_m
\]

\textbf{4. ตรวจสอบความสมเหตุสมผล}\\
ค่าจริงในช่วงนี้คือ \(V_m\) ผลประมาณสูงกว่าประมาณ 2\% เนื่องจากใช้จำนวนพจน์จำกัด

\textbf{5. แปลความหมาย}\\
ฮาร์มอนิกจำนวนเพียงสามลำดับให้ค่าภายในช่วงราบใกล้เคียงดี แต่บริเวณขอบกระโดดจะมีความคลาดเคลื่อนมากกว่า

\paragraph{ตัวอย่างที่ 10: ใช้พาร์เซวาลตรวจสอบค่า RMS ของรูปคลื่นสี่เหลี่ยม}

\textbf{ตัวอย่างเพิ่มเติมเพื่อการเรียนรู้}\\
สำหรับรูปคลื่นสี่เหลี่ยม \(\pm V_m\) ซึ่งมี \(a_0=a_n=0\) และ

\[
b_n=
\begin{cases}
\dfrac{4V_m}{n\pi}, & n\ \text{คี่}\\
0, & n\ \text{คู่}
\end{cases}
\]

จงแสดงว่า \(V_{\mathrm{rms}}=V_m\)

\textbf{1. ใช้พาร์เซวาล}

\[
V_{\mathrm{rms}}^2
=\frac{1}{2}\sum_{n=1}^{\infty}b_n^2
\]

เหลือเฉพาะ \(n\) คี่

\[
V_{\mathrm{rms}}^2
=\frac{1}{2}\frac{16V_m^2}{\pi^2}
\left(1+\frac{1}{3^2}+\frac{1}{5^2}+\cdots\right)
\]

ใช้ผลรวมอนุกรมจำนวนคี่

\[
1+\frac{1}{3^2}+\frac{1}{5^2}+\cdots=\frac{\pi^2}{8}
\]

ดังนั้น

\[
V_{\mathrm{rms}}^2
=\frac{8V_m^2}{\pi^2}\frac{\pi^2}{8}
=V_m^2
\]

จึงได้

\[
V_{\mathrm{rms}}=V_m
\]

\textbf{2. ตรวจสอบในโดเมนเวลา}\\
เพราะ \(v^2(t)=V_m^2\) ตลอดคาบ ค่าเฉลี่ยของ \(v^2(t)\) คือ \(V_m^2\) และรากที่สองคือ \(V_m\) ตรงกับผลจากพาร์เซวาล

\textbf{3. ความหมายทางกายภาพ}\\
กำลังสองเฉลี่ยทั้งหมดของรูปคลื่นเท่ากับผลรวมส่วนร่วมจากฮาร์มอนิกทุกลำดับ การตัดฮาร์มอนิกออกทำให้ค่ากำลังสองเฉลี่ยของรูปคลื่นประมาณต่ำกว่าค่าจริง

\subsubsection{7.14 สรุปขั้นตอนการแก้ปัญหาอนุกรมฟูเรียร์}

การแก้โจทย์อย่างเป็นระบบควรใช้รายการตรวจสอบต่อไปนี้

\begin{enumerate}
\def\labelenumi{\arabic{enumi}.}
\item
  \textbf{กำหนดรูปคลื่นหนึ่งคาบ} และเขียนฟังก์ชันเป็นช่วงให้ชัดเจน
\item
  \textbf{ระบุคาบและความถี่มูลฐาน} \(T\), \(f_0\), \(\omega_0\)
\item
  \textbf{เลือกจุดกำเนิดเวลา} ที่ทำให้สมมาตรเห็นชัด
\item
  \textbf{ตรวจสมมาตร} คู่ คี่ ครึ่งคลื่น หรือหนึ่งในสี่คลื่น
\item
  \textbf{หาค่าเฉลี่ย} เพื่อประเมินล่วงหน้าว่า DC ควรเป็นศูนย์หรือไม่
\item
  \textbf{เลือกสูตรสัมประสิทธิ์} ให้ตรงกับคาบ \(T\) หรือ \(2L\)
\item
  \textbf{แบ่งอินทิกรัล} ตามช่วงนิยามของฟังก์ชัน
\item
  \textbf{คำนวณและจัดรูปตามลำดับ \(n\)} ตรวจกรณี \(n\) คู่และคี่
\item
  \textbf{เขียนพจน์แรกอย่างน้อย 3--5 พจน์} เพื่อมองเห็นรูปแบบ
\item
  \textbf{ตรวจหน่วยและสมมาตรของคำตอบ}
\item
  \textbf{สร้างสเปกตรัมขนาดและเฟส} เมื่อโจทย์ต้องการตีความความถี่
\item
  \textbf{ตรวจสอบย้อนกลับ} ด้วยค่าเฉลี่ย จุดตัวอย่าง หรือพาร์เซวาลตามความเหมาะสม
\end{enumerate}

%\par\noindent\rule{\textwidth}{0.4pt}\par

\subsection{8. ความเข้าใจผิดที่พบบ่อย}

\subsubsection{8.1 พจน์ของอนุกรมเข้าใกล้ศูนย์จึงสรุปว่าอนุกรมลู่เข้า}

\textbf{ความเข้าใจที่ถูกต้อง:} \(a_n\to0\) เป็นเงื่อนไขจำเป็น แต่ไม่เพียงพอ ต้องพิจารณาผลบวกย่อย \(S_N\) ตัวอย่าง \(\sum1/n\) ลู่ออกแม้ \(1/n\to0\)

\subsubsection{8.2 ค่า \(a_0\) คือค่าเฉลี่ยของสัญญาณ}

\textbf{ความเข้าใจที่ถูกต้อง:} ภายใต้รูปแบบ

\[
f(t)=\frac{a_0}{2}+\sum\cdots
\]

ค่าเฉลี่ยคือ \(a_0/2\) ไม่ใช่ \(a_0\)

\subsubsection{8.3 เลือกช่วงอินทิเกรตไม่ครบหนึ่งคาบ}

\textbf{ความเข้าใจที่ถูกต้อง:} ช่วงอินทิเกรตเลือกเริ่มที่ใดก็ได้ แต่ความยาวต้องเท่ากับ \(T\) และนิยามฟังก์ชันต้องถูกต้องตลอดช่วง

\subsubsection{8.4 ฟังก์ชันมีพื้นที่บวกและลบเท่ากันจึงเป็นฟังก์ชันคี่}

\textbf{ความเข้าใจที่ถูกต้อง:} พื้นที่รวมเป็นศูนย์บอกเพียงค่าเฉลี่ยเป็นศูนย์ ไม่ได้ยืนยันว่า \(f(-t)=-f(t)\) ต้องตรวจสมมาตรทีละจุด

\subsubsection{8.5 รูปคลื่นสี่เหลี่ยมมีเฉพาะฮาร์มอนิกคี่ทุกกรณี}

\textbf{ความเข้าใจที่ถูกต้อง:} รูปคลื่นสี่เหลี่ยมสมมาตรครึ่งคลื่นแบบมาตรฐานมีเฉพาะฮาร์มอนิกคี่ แต่หาก duty cycle ไม่ใช่ 50\% มี DC offset หรือเลื่อนจุดอ้างอิง รูปแบบฮาร์มอนิกอาจเปลี่ยน

\subsubsection{8.6 ฟังก์ชันคู่หมายถึงมีเฉพาะฮาร์มอนิกลำดับคู่}

\textbf{ความเข้าใจที่ถูกต้อง:} คำว่า ``ฟังก์ชันคู่'' หมายถึงสมมาตร \(f(-t)=f(t)\) และทำให้พจน์ไซน์หายไป มิได้หมายความว่าฮาร์มอนิกลำดับคี่หายไป

\subsubsection{8.7 เพิ่มจำนวนพจน์แล้วปรากฏการณ์กิบส์ต้องหายไป}

\textbf{ความเข้าใจที่ถูกต้อง:} บริเวณสั่นแคบลงและความคลาดเคลื่อนโดยรวมลดลง แต่ค่าพุ่งเกินสูงสุดใกล้จุดกระโดดไม่หายทั้งหมด

\subsubsection{8.8 สเปกตรัมขนาดเพียงอย่างเดียวกำหนดรูปคลื่นได้เสมอ}

\textbf{ความเข้าใจที่ถูกต้อง:} ต้องพิจารณาเฟสด้วย เพราะเฟสกำหนดการจัดตำแหน่งขององค์ประกอบแต่ละความถี่ในโดเมนเวลา

\subsubsection{8.9 ค่าสัมประสิทธิ์ทุกตัวต้องเป็นบวก}

\textbf{ความเข้าใจที่ถูกต้อง:} \(a_n\) และ \(b_n\) อาจเป็นบวก ลบ หรือศูนย์ เครื่องหมายสะท้อนเฟสสัมพัทธ์ขององค์ประกอบนั้น

\subsubsection{8.10 จำนวนพจน์มากขึ้นย่อมเหมาะสมที่สุดในทุกงาน}

\textbf{ความเข้าใจที่ถูกต้อง:} จำนวนพจน์มากเพิ่มความแม่นยำ แต่เพิ่มภาระคำนวณและอาจเกินแบนด์วิดท์ของระบบ ต้องเลือกตามวัตถุประสงค์ ความละเอียด และข้อจำกัดจริง

%\par\noindent\rule{\textwidth}{0.4pt}\par

\subsection{9. การประยุกต์ใช้ในงานจริง}

\subsubsection{9.1 การวิเคราะห์แรงดันและกระแสที่ไม่เป็นไซน์}

โหลดอิเล็กทรอนิกส์และวงจรสวิตชิ่งอาจดึงกระแสเป็นพัลส์ แม้แรงดันแหล่งจ่ายใกล้ไซน์ อนุกรมฟูเรียร์ช่วยระบุฮาร์มอนิกของกระแสและเตรียมพื้นฐานสำหรับวิเคราะห์ผลต่อสาย ตัวนำเป็นกลาง หม้อแปลง และอุปกรณ์ชดเชย

\subsubsection{9.2 อิเล็กทรอนิกส์กำลังและ PWM}

สัญญาณ PWM ประกอบด้วยความถี่สวิตชิ่งและองค์ประกอบฮาร์มอนิก การปรับ duty cycle เปลี่ยนค่าเฉลี่ยและสเปกตรัม ดังที่เห็นจากตัวอย่างขบวนพัลส์ การออกแบบฟิลเตอร์เอาต์พุตต้องทราบว่าองค์ประกอบใดต้องผ่านและองค์ประกอบใดต้องลดทอน

\subsubsection{9.3 ระบบสื่อสาร}

รูปคลื่นดิจิทัลมีขอบชันและต้องใช้แบนด์วิดท์หลายฮาร์มอนิก หากช่องสัญญาณลดทอนความถี่สูง รูปพัลส์อาจมนและเกิดความคลาดเคลื่อนของเวลา การวิเคราะห์ฟูเรียร์ช่วยเชื่อมรูปร่างพัลส์กับแบนด์วิดท์

\subsubsection{9.4 ระบบควบคุมและการกระตุ้นเป็นคาบ}

แรงรบกวนหรือคำสั่งเป็นคาบสามารถแยกเป็นฮาร์มอนิก แต่ละฮาร์มอนิกผ่านระบบด้วยขนาดและเฟสต่างกัน ทำให้สามารถคาดการณ์การสั่นพ้องหรือการขยายบางความถี่ได้

\subsubsection{9.5 เครื่องมือวัดและการตรวจสอบสภาพ}

สัญญาณสั่นสะเทือน กระแสมอเตอร์ และแรงดันจากเซนเซอร์อาจมีองค์ประกอบความถี่สัมพันธ์กับความผิดปกติ การมองสเปกตรัมทำให้แยกองค์ประกอบที่มองไม่ชัดในกราฟเวลาได้ แม้การใช้งานจริงมักอาศัย DFT หรือ FFT แต่แนวคิดพื้นฐานมาจากอนุกรมและการวิเคราะห์ฟูเรียร์

\subsubsection{9.6 คุณภาพกำลังไฟฟ้า}

ฮาร์มอนิกเป็นตัวชี้วัดความผิดเพี้ยนของรูปคลื่น การประเมินเชิงปริมาณ เช่น RMS ของฮาร์มอนิกและ THD จะศึกษาในบทถัดไป บทนี้ให้ฐานในการระบุว่าองค์ประกอบใดเป็น DC มูลฐาน และฮาร์มอนิกลำดับสูง

\subsubsection{9.7 กรณีศึกษาเชิงแนวคิด: อินเวอร์เตอร์รูปคลื่นสี่เหลี่ยม}

อินเวอร์เตอร์ที่ให้แรงดันสี่เหลี่ยม \(\pm V_m\) มีฮาร์มอนิกคี่และขนาดลดตาม \(1/n\) หากต่อกับโหลดเชิงเหนี่ยวนำ ฮาร์มอนิกลำดับสูงจะเห็นรีแอกแตนซ์สูงขึ้นตาม \(X_L=2\pi fL\) ทำให้กระแสฮาร์มอนิกลำดับสูงลดลงมากกว่าแรงดัน จึงเป็นตัวอย่างว่ารูปคลื่นแรงดันและกระแสของโหลดอาจมีความผิดเพี้ยนไม่เท่ากัน

\subsubsection{9.8 กิจกรรมวิเคราะห์รูปคลื่นสำหรับชั้นเรียน}

\textbf{ชื่อกิจกรรม:} จากโดเมนเวลาไปสู่สเปกตรัม\\
\textbf{ผลลัพธ์ที่เกี่ยวข้อง:} CLO3, CLO5, CLO6\\
\textbf{วัตถุประสงค์:} ให้ผู้เรียนใช้คาบ สมมาตร และรูปร่างของสัญญาณทำนายลักษณะสเปกตรัมก่อนคำนวณ\\
\textbf{ระยะเวลา:} 25 นาที\\
\textbf{รูปแบบ:} กลุ่ม 3--4 คน\\
\textbf{สื่อ:} แผ่นกราฟรูปคลื่นสี่เหลี่ยม พัลส์ 25\% สามเหลี่ยม และฟันเลื่อย

\textbf{ขั้นตอน}

\begin{enumerate}
\def\labelenumi{\arabic{enumi}.}
\item
  ระบุคาบ \(T\) และความถี่มูลฐานของแต่ละรูปคลื่น
\item
  ตรวจฟังก์ชันคู่ คี่ และสมมาตรครึ่งคลื่น
\item
  ทำนายว่า DC, \(a_n\), \(b_n\), ฮาร์มอนิกคู่ หรือฮาร์มอนิกคี่ส่วนใดควรเป็นศูนย์
\item
  วาดสเปกตรัมเชิงคุณภาพโดยไม่ต้องระบุขนาดแน่นอน
\item
  เปรียบเทียบคำตอบระหว่างกลุ่มและอธิบายเหตุผล
\end{enumerate}

\textbf{ผลลัพธ์ที่คาดหวัง:} ผู้เรียนสามารถใช้สมมาตรลดการคำนวณและเชื่อมขอบชันกับการมีฮาร์มอนิกลำดับสูง

\textbf{วิธีประเมิน:} ใช้เกณฑ์ 4 ด้าน ด้านละ 2 คะแนน ได้แก่ การหาคาบ การจำแนกสมมาตร การทำนายสัมประสิทธิ์ และการอธิบายสเปกตรัม

%\par\noindent\rule{\textwidth}{0.4pt}\par

\subsection{10. สรุปท้ายบท}

\begin{enumerate}
\def\labelenumi{\arabic{enumi}.}
\item
  ลำดับคือรายการพจน์ตามดัชนี ส่วนอนุกรมคือผลบวกของพจน์ และผลบวกย่อย \(S_N\) ใช้พิจารณาการลู่เข้า
\item
  อนุกรมเรขาคณิตอนันต์ลู่เข้าเมื่อ \(|r|<1\) และมีผลบวก \(a_1/(1-r)\)
\item
  ฟังก์ชันเป็นคาบมี \(f(t+T)=f(t)\) ความถี่มูลฐานคือ \(f_0=1/T\) และ \(\omega_0=2\pi/T\)
\item
  อนุกรมฟูเรียร์ตรีโกณมิติแทนฟังก์ชันเป็นคาบด้วย DC โคไซน์ และไซน์ที่ความถี่ \(n\omega_0\)
\item
  ออร์โทกอนัลลิตีทำให้สามารถแยกสัมประสิทธิ์ของฮาร์มอนิกแต่ละลำดับได้
\item
  สำหรับคาบ \(T\)
\end{enumerate}

\[
a_0=\frac{2}{T}\int_T f(t)dt,
\quad
a_n=\frac{2}{T}\int_T f(t)\cos(n\omega_0t)dt,
\quad
b_n=\frac{2}{T}\int_T f(t)\sin(n\omega_0t)dt
\]

โดยสัญลักษณ์ \(\int_T\) หมายถึงอินทิเกรตตลอดช่วงหนึ่งคาบ

\begin{enumerate}
\def\labelenumi{\arabic{enumi}.}
\setcounter{enumi}{6}
\item
  ค่าเฉลี่ยของสัญญาณคือ \(a_0/2\)
\item
  ฟังก์ชันคู่ทำให้ \(b_n=0\) ส่วนฟังก์ชันคี่ทำให้ \(a_0=a_n=0\)
\item
  สมมาตรครึ่งคลื่นทำให้องค์ประกอบ DC และฮาร์มอนิกคู่หายไป
\item
  ขนาดฮาร์มอนิกคือ \(A_n=\sqrt{a_n^2+b_n^2}\) และเฟสตามรูป \(A_n\cos(n\omega_0t-\theta_n)\) คือ \(\theta_n=\operatorname{atan2}(b_n,a_n)\)
\item
  ที่จุดกระโดด อนุกรมฟูเรียร์ลู่เข้าหาค่าเฉลี่ยของค่าด้านซ้ายและขวา และผลรวมตัดทอนเกิดปรากฏการณ์กิบส์
\item
  ความสัมพันธ์ของพาร์เซวาลเชื่อมกำลังสองเฉลี่ยของสัญญาณกับผลรวมกำลังสองของสัมประสิทธิ์ฟูเรียร์
\item
  อนุกรมฟูเรียร์เป็นพื้นฐานของการวิเคราะห์ฮาร์มอนิก RMS THD ฟิลเตอร์ และการตอบสนองความถี่ในบทต่อไป
\end{enumerate}

%\par\noindent\rule{\textwidth}{0.4pt}\par

\subsection{11. คำถามทบทวน}

\begin{enumerate}
\def\labelenumi{\arabic{enumi}.}
\item
  ลำดับ อนุกรม และผลบวกย่อยแตกต่างกันอย่างไร
\item
  เหตุใดเงื่อนไข \(a_n\to0\) จึงไม่เพียงพอที่จะยืนยันว่า \(\sum a_n\) ลู่เข้า
\item
  คาบมูลฐานแตกต่างจากคาบที่เป็นพหุคูณของคาบมูลฐานอย่างไร
\item
  ออร์โทกอนัลลิตีของไซน์และโคไซน์ช่วยให้หาสัมประสิทธิ์ฟูเรียร์ได้อย่างไร
\item
  เหตุใดค่าเฉลี่ยของสัญญาณจึงเท่ากับ \(a_0/2\)
\item
  ฟังก์ชันคู่และฟังก์ชันคี่ส่งผลต่อ \(a_n\) และ \(b_n\) อย่างไร
\item
  สมมาตรครึ่งคลื่นกับฟังก์ชันคี่เป็นแนวคิดเดียวกันหรือไม่ จงอธิบาย
\item
  เหตุใดรูปคลื่นที่มีขอบชันจึงต้องใช้ฮาร์มอนิกลำดับสูงจำนวนมาก
\item
  ที่จุดไม่ต่อเนื่อง อนุกรมฟูเรียร์ลู่เข้าหาค่าใด
\item
  สเปกตรัมขนาดและสเปกตรัมเฟสให้ข้อมูลต่างกันอย่างไร
\item
  ปรากฏการณ์กิบส์เปลี่ยนแปลงอย่างไรเมื่อเพิ่มจำนวนพจน์
\item
  ความสัมพันธ์ของพาร์เซวาลมีความหมายต่อแรงดัน RMS อย่างไร
\end{enumerate}

%\par\noindent\rule{\textwidth}{0.4pt}\par

\subsection{12. แบบฝึกหัดท้ายบท}

\subsubsection{12.1 ระดับพื้นฐาน}

\paragraph{ข้อ 1 ลำดับและพจน์ทั่วไป}

กำหนดลำดับ \(4,7,10,13,\ldots\)

\begin{enumerate}
\def\labelenumi{\arabic{enumi}.}
\item
  จงระบุชนิดของลำดับ
\item
  จงหาพจน์ทั่วไป \(a_n\)
\item
  จงหาพจน์ที่ 25
\end{enumerate}

\paragraph{ข้อ 2 ผลบวกอนุกรมเลขคณิต}

จงหาผลบวก

\[
3+8+13+\cdots+98
\]

โดยแสดงจำนวนพจน์และขั้นตอนคำนวณ

\paragraph{ข้อ 3 อนุกรมเรขาคณิต}

พิจารณาอนุกรม

\[
12+6+3+1.5+\cdots
\]

\begin{enumerate}
\def\labelenumi{\arabic{enumi}.}
\item
  จงตรวจสอบว่าอนุกรมลู่เข้าหรือลู่ออก
\item
  ถ้าลู่เข้า จงหาผลบวกอนันต์
\item
  จงหาผลบวก 6 พจน์แรกและเปรียบเทียบกับผลบวกอนันต์
\end{enumerate}

\paragraph{ข้อ 4 คาบและฮาร์มอนิก}

สัญญาณแรงดันหนึ่งมีคาบ \(T=4\ \text{ms}\) จงหา

\begin{enumerate}
\def\labelenumi{\arabic{enumi}.}
\item
  ความถี่มูลฐาน \(f_0\)
\item
  ความถี่เชิงมุมมูลฐาน \(\omega_0\)
\item
  ความถี่ของฮาร์มอนิกลำดับที่ 2, 3 และ 7
\end{enumerate}

\paragraph{ข้อ 5 การจำแนกสมมาตร}

จงจำแนกฟังก์ชันต่อไปนี้ว่าเป็นฟังก์ชันคู่ ฟังก์ชันคี่ หรือไม่เป็นทั้งคู่และคี่

\begin{enumerate}
\def\labelenumi{\arabic{enumi}.}
\item
  \(f_1(t)=t^3\)
\item
  \(f_2(t)=t^2+2\)
\item
  \(f_3(t)=1+t\)
\item
  \(f_4(t)=\cos(3t)\)
\item
  \(f_5(t)=\sin(2t)+\cos(t)\)
\end{enumerate}

พร้อมระบุผลที่คาดต่อสัมประสิทธิ์ฟูเรียร์

\subsubsection{12.2 ระดับวิเคราะห์}

\paragraph{ข้อ 6 องค์ประกอบ DC}

ให้ฟังก์ชันคาบ \(T\) นิยามหนึ่งคาบเป็น

\[
x(t)=
\begin{cases}
6, & 0<t<T/4\\
2, & T/4<t<T
\end{cases}
\]

จงหา \(a_0\) และค่าเฉลี่ย \(x_{\mathrm{avg}}\)

\paragraph{ข้อ 7 รูปคลื่นสี่เหลี่ยมสมมาตรคี่}

รูปคลื่นแรงดัน \(\pm20\ \text{V}\) มีคาบ \(T=10\ \text{ms}\) และมีรูปเหมือนตัวอย่างที่ 5

\begin{enumerate}
\def\labelenumi{\arabic{enumi}.}
\item
  จงเขียนอนุกรมฟูเรียร์สามพจน์ที่ไม่เป็นศูนย์แรก
\item
  จงระบุความถี่และแอมพลิจูดของฮาร์มอนิกลำดับที่ 1, 3 และ 5
\item
  อธิบายเหตุผลที่ฮาร์มอนิกคู่เป็นศูนย์
\end{enumerate}

\paragraph{ข้อ 8 พัลส์ duty cycle 25\%}

ขบวนพัลส์สมมาตรคู่มีแอมพลิจูด \(A=8\ \text{V}\) คาบ \(T=1\ \text{ms}\) และ duty cycle \(D=0.25\)

\begin{enumerate}
\def\labelenumi{\arabic{enumi}.}
\item
  จงหาค่า DC
\item
  จงหา \(a_1\), \(a_2\), \(a_3\) และ \(a_4\)
\item
  จงระบุว่าฮาร์มอนิกลำดับใดในชุดนี้เป็นศูนย์
\end{enumerate}

ใช้

\[
a_n=\frac{2A}{n\pi}\sin(n\pi D)
\]

\paragraph{ข้อ 9 การสังเคราะห์จากสัมประสิทธิ์}

กำหนด

\[
x(t)=2+3\cos(100\pi t)+4\sin(100\pi t)-2\cos(200\pi t)
\]

\begin{enumerate}
\def\labelenumi{\arabic{enumi}.}
\item
  จงหา \(T\) และ \(f_0\)
\item
  จงระบุ \(a_0\), \(a_1\), \(b_1\), \(a_2\), \(b_2\)
\item
  จงแปลงฮาร์มอนิกลำดับที่ 1 เป็นรูปขนาด--เฟส
\item
  จงวาดสเปกตรัมขนาดเชิงเส้น
\end{enumerate}

\paragraph{ข้อ 10 การลู่เข้าที่จุดกระโดด}

ฟังก์ชันคาบหนึ่งกระโดดจาก \(-3\ \text{V}\) ทางซ้ายของ \(t=t_d\) ไปเป็น \(7\ \text{V}\) ทางขวา

\begin{enumerate}
\def\labelenumi{\arabic{enumi}.}
\item
  อนุกรมฟูเรียร์ลู่เข้าหาค่าใดที่ \(t=t_d\)
\item
  ขนาดการกระโดดเท่ากับเท่าใด
\item
  ประมาณค่าพุ่งเกินเหนือระดับบนจากปรากฏการณ์กิบส์ โดยใช้ร้อยละ 9 ของขนาดการกระโดด
\item
  เมื่อเพิ่มจำนวนฮาร์มอนิก ค่าพุ่งเกินและความกว้างบริเวณสั่นเปลี่ยนอย่างไร
\end{enumerate}

\subsubsection{12.3 ระดับประยุกต์ใช้}

\paragraph{ข้อ 11 การวิเคราะห์สัญญาณ PWM สำหรับโหลด DC}

วงจรสวิตชิ่งสร้างพัลส์แรงดัน \(0\) ถึง \(24\ \text{V}\) ที่ความถี่สวิตชิ่ง 2 kHz และ duty cycle 40\% โดยสมมติพัลส์ถูกจัดให้อยู่กึ่งกลางคาบ

\begin{enumerate}
\def\labelenumi{\arabic{enumi}.}
\item
  จงหาคาบสวิตชิ่ง
\item
  จงหาค่าเฉลี่ยของแรงดัน
\item
  จงหาสัมประสิทธิ์โคไซน์ \(a_1\), \(a_2\) และ \(a_5\)
\item
  อธิบายว่าทำไมตัวเหนี่ยวนำอนุกรมจึงช่วยให้กระแสเรียบขึ้น โดยเชื่อมโยงกับความถี่ฮาร์มอนิก
\end{enumerate}

\paragraph{ข้อ 12 การเลือกจำนวนฮาร์มอนิกสำหรับการสื่อสารพัลส์}

ต้องส่งสัญญาณสี่เหลี่ยมสมมาตรคี่ที่ \(f_0=1\ \text{kHz}\) ผ่านช่องสัญญาณที่ส่งผ่านได้ถึง 7 kHz และลดทอนความถี่สูงกว่านี้อย่างมาก

\begin{enumerate}
\def\labelenumi{\arabic{enumi}.}
\item
  ฮาร์มอนิกใดของสัญญาณสี่เหลี่ยมผ่านช่องสัญญาณได้
\item
  จงเขียนรูปคลื่นประมาณที่ปลายทางโดยสมมติฮาร์มอนิกที่ผ่านไม่ถูกลดทอนและไม่เปลี่ยนเฟส
\item
  เปรียบเทียบกับสัญญาณต้นฉบับในด้านขอบชันและการสั่น
\item
  อธิบายว่าการเพิ่มแบนด์วิดท์ถึง 15 kHz จะเปลี่ยนรูปคลื่นอย่างไร
\end{enumerate}

\paragraph{ข้อ 13 พาร์เซวาลและสัดส่วน RMS จากฮาร์มอนิก}

แรงดันเป็นคาบมีองค์ประกอบ

\[
v(t)=10+20\cos(\omega_0t)+6\sin(3\omega_0t)\ \text{V}
\]

\begin{enumerate}
\def\labelenumi{\arabic{enumi}.}
\item
  จงหาค่า \(V_{\mathrm{rms}}\) ด้วยความสัมพันธ์ของพาร์เซวาล
\item
  จงหาส่วนร่วมกำลังสองเฉลี่ยจาก DC มูลฐาน และฮาร์มอนิกลำดับที่ 3
\item
  ถ้าฟิลเตอร์ตัดฮาร์มอนิกลำดับที่ 3 ออกโดยไม่เปลี่ยนส่วนอื่น ค่า RMS ใหม่เป็นเท่าใด
\item
  อธิบายความแตกต่างของรูปคลื่นก่อนและหลังฟิลเตอร์เชิงคุณภาพ
\end{enumerate}

%\par\noindent\rule{\textwidth}{0.4pt}\par

\subsubsection{12.4 แนวคำตอบและเฉลยสำหรับผู้สอน}

\begin{quote}
\textbf{หมายเหตุ:} ส่วนนี้ควรแยกออกจากเอกสารสำหรับนักศึกษาเมื่อต้องการใช้แบบฝึกหัดเป็นงานส่งหรือการประเมิน
\end{quote}

\paragraph{เฉลยข้อ 1}

\textbf{แนวคิด:} ลำดับเลขคณิตมีผลต่างร่วมคงที่

\[
d=7-4=3
\]

พจน์ทั่วไป

\[
a_n=a_1+(n-1)d
\]

\[
a_n=4+3(n-1)=3n+1
\]

พจน์ที่ 25

\[
a_{25}=3(25)+1=76
\]

\textbf{คำตอบ:} เป็นลำดับเลขคณิต \(a_n=3n+1\) และ \(a_{25}=76\)

\textbf{ข้อผิดพลาดที่พบบ่อย:} ใช้ \(a_n=a_1+nd\) ทำให้ดัชนีคลาดหนึ่งตำแหน่ง

\paragraph{เฉลยข้อ 2}

\textbf{แนวคิด:} ต้องหาจำนวนพจน์ก่อน

\[
a_N=a_1+(N-1)d
\]

โดย \(a_1=3\), \(d=5\), \(a_N=98\)

\[
98=3+(N-1)5
\]

\[
95=5(N-1)
\]

\[
N=20
\]

ผลบวก

\[
S_{20}=\frac{20}{2}(3+98)=10(101)=1010
\]

\textbf{คำตอบ:} มี 20 พจน์ และผลบวกเท่ากับ \(1010\)

\paragraph{เฉลยข้อ 3}

\textbf{แนวคิด:} เป็นอนุกรมเรขาคณิต \(a_1=12\), \(r=0.5\)

เพราะ

\[
|r|=0.5<1
\]

จึงลู่เข้า

ผลบวกอนันต์

\[
S_\infty=\frac{12}{1-0.5}=24
\]

ผลบวก 6 พจน์

\[
S_6=12\frac{1-0.5^6}{1-0.5}
\]

\[
S_6=24\left(1-\frac{1}{64}\right)=23.625
\]

ความแตกต่างจากผลบวกอนันต์

\[
24-23.625=0.375
\]

\textbf{คำตอบ:} อนุกรมลู่เข้าสู่ 24; หกพจน์แรกมีผลบวก 23.625 ซึ่งต่ำกว่าค่าลิมิต 0.375

\paragraph{เฉลยข้อ 4}

แปลงคาบ

\[
T=4\times10^{-3}\ \text{s}
\]

ความถี่มูลฐาน

\[
f_0=\frac{1}{T}=250\ \text{Hz}
\]

ความถี่เชิงมุม

\[
\omega_0=2\pi(250)=500\pi\approx1570.8\ \text{rad/s}
\]

ฮาร์มอนิก

\[
f_2=500\ \text{Hz},\quad f_3=750\ \text{Hz},\quad f_7=1750\ \text{Hz}
\]

\paragraph{เฉลยข้อ 5}

\begin{enumerate}
\def\labelenumi{\arabic{enumi}.}
\item
  \(t^3\) เป็นคี่ จึง \(a_0=a_n=0\)
\item
  \(t^2+2\) เป็นคู่ จึง \(b_n=0\)
\item
  \(1+t\) ไม่เป็นทั้งคู่และคี่ โดยทั่วไปมี DC โคไซน์ และไซน์
\item
  \(\cos(3t)\) เป็นคู่ จึงมีเฉพาะโคไซน์ที่สัมพันธ์กับคาบที่เลือก
\item
  \(\sin(2t)+\cos(t)\) ไม่เป็นทั้งคู่และคี่ เพราะเป็นผลรวมของส่วนคี่และส่วนคู่
\end{enumerate}

\textbf{เหตุผล:} ต้องแทน \(-t\) แล้วเปรียบเทียบกับ \(f(t)\) และ \(-f(t)\)

\paragraph{เฉลยข้อ 6}

\textbf{แนวคิด:} หาค่าเฉลี่ยจากพื้นที่ในหนึ่งคาบ

\[
x_{\mathrm{avg}}=\frac{1}{T}\left[
\int_0^{T/4}6\,dt+\int_{T/4}^{T}2\,dt
\right]
\]

\[
x_{\mathrm{avg}}=\frac{1}{T}\left[6\frac{T}{4}+2\frac{3T}{4}\right]
\]

\[
x_{\mathrm{avg}}=\frac{1}{T}\left(\frac{6T}{4}+\frac{6T}{4}\right)=3
\]

ดังนั้น

\[
a_0=2x_{\mathrm{avg}}=6
\]

\textbf{คำตอบ:} \(a_0=6\) และ \(a_0/2=x_{\mathrm{avg}}=3\)

\paragraph{เฉลยข้อ 7}

ให้ \(V_m=20\ \text{V}\) และ

\[
f_0=\frac{1}{10\times10^{-3}}=100\ \text{Hz}
\]

\[
\omega_0=200\pi\ \text{rad/s}
\]

อนุกรมสามพจน์แรก

\[
v(t)\approx\frac{80}{\pi}\left[
\sin(200\pi t)+\frac{1}{3}\sin(600\pi t)+\frac{1}{5}\sin(1000\pi t)
\right]\ \text{V}
\]

ขนาดฮาร์มอนิก

\[
A_1=\frac{80}{\pi}\approx25.465\ \text{V}
\]

\[
A_3=\frac{80}{3\pi}\approx8.488\ \text{V}
\]

\[
A_5=\frac{80}{5\pi}\approx5.093\ \text{V}
\]

ความถี่คือ 100, 300 และ 500 Hz

\textbf{เหตุผลที่ฮาร์มอนิกคู่เป็นศูนย์:} รูปคลื่นมีสมมาตรครึ่งคลื่น ทำให้ \(1-(-1)^n\) เป็นศูนย์เมื่อ \(n\) คู่

\paragraph{เฉลยข้อ 8}

ให้ \(A=8\), \(D=0.25\)

ค่า DC

\[
AD=8(0.25)=2\ \text{V}
\]

สัมประสิทธิ์

\[
a_n=\frac{16}{n\pi}\sin\left(\frac{n\pi}{4}\right)
\]

ลำดับที่ 1

\[
a_1=\frac{16}{\pi}\frac{\sqrt2}{2}
=\frac{8\sqrt2}{\pi}\approx3.601\ \text{V}
\]

ลำดับที่ 2

\[
a_2=\frac{16}{2\pi}\sin\left(\frac{\pi}{2}\right)
=\frac{8}{\pi}\approx2.546\ \text{V}
\]

ลำดับที่ 3

\[
a_3=\frac{16}{3\pi}\sin\left(\frac{3\pi}{4}\right)
=\frac{8\sqrt2}{3\pi}\approx1.200\ \text{V}
\]

ลำดับที่ 4

\[
a_4=\frac{16}{4\pi}\sin(\pi)=0
\]

\textbf{คำตอบ:} DC = 2 V และฮาร์มอนิกลำดับที่ 4 เป็นศูนย์ ในภาพรวมทุกลำดับที่เป็นพหุคูณของ 4 จะเป็นศูนย์สำหรับ \(D=0.25\)

\paragraph{เฉลยข้อ 9}

จาก \(100\pi=2\pi f_0\)

\[
f_0=50\ \text{Hz},\qquad T=0.02\ \text{s}
\]

เปรียบเทียบกับรูปมาตรฐาน

\[
\frac{a_0}{2}=2\Rightarrow a_0=4
\]

\[
a_1=3,\quad b_1=4,\quad a_2=-2,\quad b_2=0
\]

ฮาร์มอนิกที่ 1

\[
A_1=\sqrt{3^2+4^2}=5
\]

\[
\theta_1=\operatorname{atan2}(4,3)=53.13^\circ
\]

ดังนั้น

\[
3\cos(100\pi t)+4\sin(100\pi t)
=5\cos(100\pi t-53.13^\circ)
\]

สเปกตรัมขนาดแบบด้านเดียวมี

\begin{itemize}
\item
  0 Hz: 2
\item
  50 Hz: 5
\item
  100 Hz: 2
\end{itemize}

เครื่องหมายลบของ \(a_2\) สะท้อนเฟส 180 องศา แต่ขนาดเท่ากับ 2

\paragraph{เฉลยข้อ 10}

ค่าที่จุดกระโดด

\[
S(t_d)=\frac{-3+7}{2}=2\ \text{V}
\]

ขนาดการกระโดด

\[
\Delta=7-(-3)=10\ \text{V}
\]

ค่าพุ่งเกินเหนือระดับบนโดยประมาณ

\[
0.09\Delta=0.9\ \text{V}
\]

ดังนั้นค่าสูงสุดโดยประมาณใกล้จุดกระโดดคือ

\[
7+0.9=7.9\ \text{V}
\]

เมื่อเพิ่มฮาร์มอนิก บริเวณสั่นจะแคบลงและเข้าใกล้จุดกระโดด แต่ค่าพุ่งเกินสัดส่วนประมาณเดิมไม่หายทั้งหมด

\paragraph{เฉลยข้อ 11}

\textbf{1. คาบสวิตชิ่ง}

\[
T=\frac{1}{2000}=0.0005\ \text{s}=0.5\ \text{ms}
\]

\textbf{2. ค่าเฉลี่ย}

\[
V_{\mathrm{avg}}=AD=24(0.4)=9.6\ \text{V}
\]

\textbf{3. สัมประสิทธิ์}

สำหรับพัลส์สมมาตรคู่

\[
a_n=\frac{2A}{n\pi}\sin(n\pi D)
=\frac{48}{n\pi}\sin(0.4n\pi)
\]

ลำดับที่ 1

\[
a_1=\frac{48}{\pi}\sin(0.4\pi)
\approx14.531\ \text{V}
\]

ลำดับที่ 2

\[
a_2=\frac{24}{\pi}\sin(0.8\pi)
\approx4.490\ \text{V}
\]

ลำดับที่ 5

\[
a_5=\frac{48}{5\pi}\sin(2\pi)=0
\]

\textbf{4. การตีความตัวเหนี่ยวนำ}

ฮาร์มอนิกลำดับที่ \(n\) มีความถี่ \(nf_0\) และเห็นรีแอกแตนซ์

\[
X_{L,n}=2\pi(nf_0)L=nX_{L,1}
\]

จึงต้านการเปลี่ยนแปลงกระแสที่ความถี่สูงมากกว่า ทำให้ริปเปิลฮาร์มอนิกลดลงและกระแสเรียบขึ้น ส่วน DC ไม่ถูกต้านด้วยรีแอกแตนซ์เชิงอุดมคติของตัวเหนี่ยวนำ

\paragraph{เฉลยข้อ 12}

สัญญาณสี่เหลี่ยมสมมาตรคี่มีฮาร์มอนิกคี่ \(f_0,3f_0,5f_0,7f_0,\ldots\)

เมื่อ \(f_0=1\) kHz และแบนด์วิดท์ถึง 7 kHz ฮาร์มอนิกที่ผ่านคือ

\[
1,3,5,7\ \text{kHz}
\]

รูปคลื่นประมาณ

\[
x_{\mathrm{out}}(t)\approx\frac{4A}{\pi}\left[
\sin(2\pi1000t)
+\frac{1}{3}\sin(2\pi3000t)
+\frac{1}{5}\sin(2\pi5000t)
+\frac{1}{7}\sin(2\pi7000t)
\right]
\]

รูปคลื่นปลายทางมีขอบมนกว่าต้นฉบับและมีการสั่นใกล้จุดเปลี่ยน เพราะฮาร์มอนิกสูงกว่าถูกตัดออก

ถ้าแบนด์วิดท์เพิ่มเป็น 15 kHz จะผ่านฮาร์มอนิกคี่ถึงลำดับ 15 ได้แก่ 1,3,5,7,9,11,13,15 kHz ทำให้ขอบชันขึ้นและช่วงเปลี่ยนแคบลง แต่ยังพบกิบส์ใกล้จุดกระโดด

\paragraph{เฉลยข้อ 13}

กำหนด

\[
\frac{a_0}{2}=10\Rightarrow a_0=20
\]

มูลฐานมี \(a_1=20\), \(b_1=0\) และฮาร์มอนิกที่ 3 มี \(a_3=0\), \(b_3=6\)

ใช้พาร์เซวาล

\[
V_{\mathrm{rms}}^2
=\frac{a_0^2}{4}+\frac{1}{2}\sum(a_n^2+b_n^2)
\]

\[
V_{\mathrm{rms}}^2
=100+\frac{1}{2}(20^2)+\frac{1}{2}(6^2)
\]

\[
V_{\mathrm{rms}}^2=100+200+18=318
\]

\[
V_{\mathrm{rms}}=\sqrt{318}\approx17.833\ \text{V}
\]

ส่วนร่วมกำลังสองเฉลี่ย

\begin{itemize}
\item
  DC: \(100\ \text{V}^2\)
\item
  มูลฐาน: \(200\ \text{V}^2\)
\item
  ฮาร์มอนิกที่ 3: \(18\ \text{V}^2\)
\end{itemize}

เมื่อตัดฮาร์มอนิกที่ 3

\[
V_{\mathrm{rms,new}}^2=100+200=300
\]

\[
V_{\mathrm{rms,new}}=\sqrt{300}\approx17.321\ \text{V}
\]

รูปคลื่นหลังฟิลเตอร์จะเรียบและใกล้ผลรวมของ DC กับไซน์มูลฐานมากขึ้น ความกระเพื่อมที่เกิดจากความถี่ \(3f_0\) หายไป และค่า RMS ลดลงเล็กน้อย

%\par\noindent\rule{\textwidth}{0.4pt}\par

\subsection{13. ข้อเสนอแนะสำหรับผู้สอนและแนวทางประเมินผล}

\subsubsection{13.1 แผนการสอนที่แนะนำ}

\paragraph{ครั้งที่ 1: ลำดับ อนุกรม และฟังก์ชันเป็นคาบ --- 3 ชั่วโมง}

\begin{table}[htbp]
    \centering
    \begin{tabular}{rlll}
        \hline
        \textbf{ช่วงเวลา} & \textbf{เนื้อหา/กิจกรรม} & \textbf{วิธีดำเนินการ} & \textbf{หลักฐานระหว่างเรียน} \\
        \hline
        0--20 นาที & ทบทวนตรีโกณมิติ อินทิกรัล คาบและความถี่ & คำถามวินิจฉัย 5 ข้อ & คำตอบรายบุคคล \\
        20--60 นาที & ลำดับ อนุกรม ผลบวกย่อย และการลู่เข้า & บรรยายสลับตัวอย่าง 1--2 & วิธีทำบนใบงาน \\
        60--90 นาที & อนุกรมเลขคณิตและเรขาคณิต & คำนวณเป็นคู่และเปรียบเทียบคำตอบ & ใบงานคู่ \\
        90--105 นาที & พัก & --- & --- \\
        105--140 นาที & ฟังก์ชันเป็นคาบ \(T,f_0,\omega_0\) & วิเคราะห์กราฟรูปคลื่นไฟฟ้า & ตารางคาบและฮาร์มอนิก \\
        140--170 นาที & แนวคิดการสังเคราะห์ฟูเรียร์ & ใช้ภาพรูปที่ 5 และอภิปราย & คำอธิบายเชิงแนวคิด \\
        170--180 นาที & Exit ticket & 3 คำถามสั้น & แบบประเมินท้ายคาบ \\
        \hline
    \end{tabular}
\end{table}

\paragraph{ครั้งที่ 2: สัมประสิทธิ์ สมมาตร และสเปกตรัม --- 3 ชั่วโมง}

\begin{table}[htbp]
    \centering
    \begin{tabular}{rlll}
        \hline
        \textbf{ช่วงเวลา} & \textbf{เนื้อหา/กิจกรรม} & \textbf{วิธีดำเนินการ} & \textbf{หลักฐานระหว่างเรียน} \\
        \hline
        0--25 นาที & ออร์โทกอนัลลิตี & วิเคราะห์พื้นที่บวก--ลบจากกราฟ & คำตอบอธิบาย \\
        25--70 นาที & ที่มาของ \(a_0,a_n,b_n\) & บรรยายแบบถามตอบ & สมการที่ผู้เรียนเติม \\
        70--105 นาที & ฟังก์ชันคู่ คี่ และครึ่งคลื่น & กิจกรรมจำแนกสมมาตร & ตารางสรุปสัมประสิทธิ์ \\
        105--120 นาที & พัก & --- & --- \\
        120--155 นาที & ตัวอย่างรูปคลื่นสี่เหลี่ยมและพัลส์ & แก้โจทย์เป็นขั้นตอน & วิธีทำกลุ่ม \\
        155--170 นาที & สเปกตรัม กิบส์ และพาร์เซวาล & วิเคราะห์ภาพและกรณีศึกษา & แผนภาพสเปกตรัม \\
        170--180 นาที & แบบทดสอบสั้น & คำนวณ 1 ข้อ + อธิบาย 1 ข้อ & กระดาษคำตอบ \\
        \hline
    \end{tabular}
\end{table}

\subsubsection{13.2 จุดที่ผู้เรียนมักมีปัญหา}

\begin{enumerate}
\def\labelenumi{\arabic{enumi}.}
\item
  สับสนระหว่างดัชนี \(n\) กับเวลา \(t\)
\item
  ไม่แยกพจน์ \(a_0\) ออกจากค่าเฉลี่ย \(a_0/2\)
\item
  เลือก \(\omega_0\) ผิดจาก \(2\pi/T\)
\item
  ใช้สูตรช่วง \([-L,L]\) ปะปนกับสูตรคาบ \(T\) โดยไม่แทน \(T=2L\)
\item
  มองสมมาตรจากรูปร่างโดยไม่ตรวจสมการ \(f(-t)\)
\item
  ลืมแบ่งอินทิกรัลเมื่อฟังก์ชันนิยามเป็นช่วง
\item
  จัดรูป \(1-(-1)^n\) ไม่ถูกต้อง
\item
  แปลงจาก \(a_n,b_n\) เป็นขนาด--เฟสโดยใช้เครื่องหมายเฟสผิด
\item
  คาดหวังให้ผลรวมตัดทอนตรงรูปคลื่นที่จุดกระโดด
\item
  วาดสเปกตรัมโดยวางแท่งที่ความถี่ต่อเนื่องแทนที่จะเป็น \(nf_0\)
\end{enumerate}

\subsubsection{13.3 วิธีช่วยผู้เรียนที่มีพื้นฐานต่างกัน}

\begin{itemize}
\item
  \textbf{ผู้เรียนที่ต้องการเสริมพื้นฐาน:} ให้ตารางอินทิกรัลและแบบฝึกแยกสมมาตรก่อนโจทย์เต็ม
\item
  \textbf{ผู้เรียนระดับปานกลาง:} ให้โครงขั้นตอนโดยเว้นช่องสำหรับเลือกช่วงและสูตร
\item
  \textbf{ผู้เรียนที่พร้อมสูง:} ให้โจทย์ย้อนกลับจากสเปกตรัมไปสร้างรูปคลื่น และให้เปรียบเทียบรูปตรีโกณมิติกับรูปเชิงซ้อน
\item
  \textbf{การเรียนแบบกลุ่ม:} กำหนดบทบาทผู้วิเคราะห์สมมาตร ผู้คำนวณอินทิกรัล ผู้ตรวจหน่วย และผู้อธิบายสเปกตรัม
\end{itemize}

\subsubsection{13.4 แนวทางประเมินระหว่างเรียน}

\begin{table}[htbp]
    \centering
    \begin{tabular}{lrll}
        \hline
        \textbf{รายการประเมิน} & \textbf{สัดส่วนภายในบท} & \textbf{CLO ที่ประเมิน} & \textbf{เกณฑ์หลัก} \\
        \hline
        แบบทดสอบวินิจฉัยและ Exit ticket & 10\% & CLO1, CLO3 & นิยาม คาบ ความถี่ และการตีความ \\
        ใบงานอนุกรมและผลบวกย่อย & 15\% & CLO1, CLO2 & สูตร ขั้นตอน และความถูกต้อง \\
        ใบงานสมมาตรและสัมประสิทธิ์ & 20\% & CLO4, CLO5 & การเลือกสูตรและลดรูป \\
        กิจกรรมสเปกตรัมกลุ่ม & 15\% & CLO3, CLO6 & การเชื่อมโดเมนเวลา--ความถี่ \\
        แบบฝึกหัดท้ายบท & 25\% & CLO1--CLO6 & การคำนวณและประยุกต์ \\
        แบบทดสอบสรุปบท & 15\% & CLO3--CLO6 & โจทย์คำนวณและคำอธิบาย \\
        \hline
    \end{tabular}
\end{table}

\begin{quote}
สัดส่วนในตารางเป็นข้อเสนอสำหรับการประเมินภายในบท ผู้สอนควรปรับให้สอดคล้องกับแผนประเมินผลของรายวิชา
\end{quote}

\subsubsection{13.5 รูบริกย่อสำหรับโจทย์หาสัมประสิทธิ์ฟูเรียร์ 10 คะแนน}

\begin{table}[htbp]
    \centering
    \begin{tabular}{llll}
        \hline
        \textbf{เกณฑ์} & \textbf{2 คะแนน} & \textbf{1 คะแนน} & \textbf{0 คะแนน} \\
        \hline
        กำหนดคาบและ \(\omega_0\) & ถูกต้องครบ & ผิดเล็กน้อยแต่แนวคิดถูก & ไม่ระบุหรือผิดหลัก \\
        วิเคราะห์สมมาตร & ระบุพร้อมเหตุผล & ระบุถูกแต่ไม่ให้เหตุผล & ระบุผิด \\
        ตั้งอินทิกรัล & ช่วงและฟังก์ชันถูก & ผิดเล็กน้อยหนึ่งจุด & ใช้สูตรไม่สอดคล้อง \\
        คำนวณสัมประสิทธิ์ & ขั้นตอนและผลถูก & มีข้อผิดพลาดคำนวณ & ไม่สามารถคำนวณได้ \\
        ตีความอนุกรม/สเปกตรัม & ระบุ DC ฮาร์มอนิก หน่วยครบ & ตีความได้บางส่วน & ไม่ตีความ \\
        \hline
    \end{tabular}
\end{table}

\subsubsection{13.6 เกณฑ์ตรวจหลักฐานการเรียนรู้}

ผู้เรียนถือว่าบรรลุผลลัพธ์ขั้นต่ำเมื่อ

\begin{itemize}
\item
  อธิบายความแตกต่างระหว่างลำดับและอนุกรมได้ถูกต้องอย่างน้อย 70\%
\item
  คำนวณอนุกรมเลขคณิตและเรขาคณิตได้อย่างน้อย 70\%
\item
  หา \(T,f_0,\omega_0\) และฮาร์มอนิกได้ถูกต้องอย่างน้อย 80\%
\item
  ตั้งสูตรสัมประสิทธิ์ฟูเรียร์และเลือกช่วงได้ถูกต้องอย่างน้อย 70\%
\item
  ใช้สมมาตรเพื่อตัดสัมประสิทธิ์ที่เป็นศูนย์ได้อย่างน้อย 70\%
\item
  อ่านสเปกตรัมและอธิบาย DC มูลฐาน และฮาร์มอนิกได้อย่างน้อย 70\%
\end{itemize}

%\par\noindent\rule{\textwidth}{0.4pt}\par

\subsection{14. สื่อ แหล่งเรียนรู้เพิ่มเติม และเอกสารอ้างอิง}

\subsubsection{14.1 สื่อแนะนำสำหรับผู้เรียน}

\begin{table}[htbp]
    \centering
    \begin{tabular}{lllll}
        \hline
        \textbf{องค์กร/ผู้จัดทำ} & \textbf{หัวข้อหรือหน้าสื่อที่ตรวจสอบได้} & \textbf{คำค้นแนะนำ} & \textbf{จุดประสงค์ในการใช้} & \textbf{ช่วงเวลาที่เหมาะสม} \\
        \hline
        MIT OpenCourseWare & Fourier Series: Basics, วิชา Differential Equations & \texttt{MIT\ OCW\ Fourier\ Series\ Basics\ definition\ coefficients} & ทบทวนนิยาม สูตรสัมประสิทธิ์ และฟังก์ชันคาบ & หลังหัวข้อ 7.5--7.7 \\
        MIT OpenCourseWare & Operations on Fourier Series & \texttt{MIT\ OCW\ operations\ Fourier\ series\ even\ odd\ convergence} & ศึกษาสมมาตร การลู่เข้า และกิบส์ & หลังหัวข้อ 7.9 และ 7.12 \\
        MIT OpenCourseWare, Signals and Systems & Continuous-Time Fourier Series & \texttt{MIT\ RES\ 6-007\ continuous\ time\ Fourier\ series} & มองอนุกรมฟูเรียร์จากมุมสัญญาณและระบบ & ก่อนเข้าสู่หัวข้อสเปกตรัม \\
        Stanford Engineering Everywhere & EE261: The Fourier Transform and Its Applications & \texttt{Stanford\ EE261\ Fourier\ series\ Brad\ Osgood} & ขยายความเข้าใจเรื่องฐานฟังก์ชัน สเปกตรัม และการประยุกต์ & สำหรับผู้เรียนที่ต้องการเนื้อหาเพิ่มเติม \\
        NPTEL, IIT Roorkee & Mathematical Methods and its Applications: Fourier series and convergence & \texttt{NPTEL\ Fourier\ series\ convergence\ IIT\ Roorkee} & ทบทวนการลู่เข้าและตัวอย่างคำนวณระดับวิศวกรรม & หลังเรียนจบบท \\
        NPTEL, IIT Kanpur & Principles of Signals and Systems & \texttt{NPTEL\ principles\ signals\ systems\ Fourier\ series} & เชื่อมโยงสัญญาณเป็นคาบกับระบบไฟฟ้า & ก่อนเริ่มบทที่ 5 \\
        Rice University ECE & Fundamentals of Electrical Engineering I: Frequency Domain & \texttt{Rice\ ECE\ classic\ Fourier\ series\ electrical\ engineering} & อ่านตัวอย่างสัญญาณและสเปกตรัมในบริบทวิศวกรรมไฟฟ้า & ใช้ประกอบกรณีศึกษา \\
        \hline
    \end{tabular}
\end{table}

\subsubsection{14.2 แนวทางใช้สื่ออย่างมีประสิทธิภาพ}

\begin{enumerate}
\def\labelenumi{\arabic{enumi}.}
\item
  ให้ผู้เรียนอ่านหรือรับชมโดยมีคำถามนำ ไม่ควรมอบสื่อโดยไม่มีภารกิจ
\item
  คำถามนำที่แนะนำ ได้แก่ ``สูตรสัมประสิทธิ์เกิดจากสมบัติใด'' ``ที่จุดกระโดดผลรวมลู่เข้าค่าใด'' และ ``สมมาตรใดทำให้ฮาร์มอนิกคู่หายไป''
\item
  หลังรับชม ให้ผู้เรียนเขียนสรุปหนึ่งย่อหน้าและแก้โจทย์สั้นหนึ่งข้อ
\item
  เปรียบเทียบรูปแบบสัญลักษณ์ของแต่ละแหล่ง เพราะบางแหล่งใช้ \(A_0\), \(a_0\) หรือกำหนดพจน์คงที่ต่างรูปกัน
\item
  เน้นให้ผู้เรียนตรวจสูตรกับนิยามของแหล่งนั้นก่อนแทนค่า เพื่อป้องกันความคลาดเคลื่อนด้วยตัวประกอบ 2
\end{enumerate}

%\par\noindent\rule{\textwidth}{0.4pt}\par

\subsubsection{14.3 เอกสารอ้างอิง}

\paragraph{14.3.1 เอกสารหลักที่ใช้กำหนดโครงสร้าง}

\begin{itemize}
\item
  \textit{teaching\_document\_structure\_5571102\_electrical\_engineering\_mathematics.md}. (2569). โครงสร้างเอกสารประกอบการสอนรายวิชา 5571102 คณิตศาสตร์วิศวกรรมไฟฟ้า.
\item
  \textit{prompt-core.md}. (2569). Prompt กลางสำหรับออกแบบเอกสารประกอบการสอนระดับอุดมศึกษา.
\item
  \textit{module-a-images.md}. (2569). Module A: ภาพประกอบและ Prompt สร้างภาพ.
\item
  \textit{module-b-equations.md}. (2569). Module B: สมการและตัวอย่างคำนวณ.
\item
  \textit{module-e-exercises-solutions.md}. (2569). Module E: แบบฝึกหัดและเฉลย.
\item
  \textit{module-f-obe-assessment.md}. (2569). Module F: OBE และการประเมินผล.
\item
  \textit{module-h-media-references.md}. (2569). Module H: สื่อ วิดีโอ และเอกสารอ้างอิง.
\end{itemize}

\paragraph{14.3.2 หนังสือและตำรา}

Brown, J. W., \& Churchill, R. V. (2011). \textit{Fourier series and boundary value problems} (8th ed.). McGraw-Hill.

Johnson, D. H. (2016). \textit{Fundamentals of electrical engineering I}. Rice University.

Kreyszig, E. (2011). \textit{Advanced engineering mathematics} (10th ed.). John Wiley \& Sons.

Oppenheim, A. V., Willsky, A. S., \& Nawab, S. H. (1997). \textit{Signals and systems} (2nd ed.). Prentice Hall.

Osgood, B. (2007). \textit{The Fourier transform and its applications: Lecture notes for EE 261}. Stanford University.

\paragraph{14.3.3 แหล่งข้อมูลออนไลน์จากมหาวิทยาลัยและหน่วยงานการศึกษา}

Massachusetts Institute of Technology OpenCourseWare. (2011). \textit{Fourier series: Basics}. In \textit{18.03SC Differential Equations}. \url{https://ocw.mit.edu/courses/18-03sc-differential-equations-fall-2011/pages/unit-iii-fourier-series-and-laplace-transform/fourier-series-basics/}

Massachusetts Institute of Technology OpenCourseWare. (2011). \textit{Operations on Fourier series}. In \textit{18.03SC Differential Equations}. \url{https://ocw.mit.edu/courses/18-03sc-differential-equations-fall-2011/pages/unit-iii-fourier-series-and-laplace-transform/operations-on-fourier-series/}

Massachusetts Institute of Technology OpenCourseWare. (2011). \textit{Continuous-time Fourier series}. In \textit{RES.6-007 Signals and Systems}. \url{https://ocw.mit.edu/courses/res-6-007-signals-and-systems-spring-2011/resources/lecture-7-continuous-time-fourier-series/}

National Programme on Technology Enhanced Learning. (n.d.). \textit{Mathematical methods and its applications}. Indian Institute of Technology Roorkee. \url{https://nptel.ac.in/courses/111107098}

National Programme on Technology Enhanced Learning. (n.d.). \textit{Principles of signals and systems}. Indian Institute of Technology Kanpur. \url{https://nptel.ac.in/courses/108104100}

Rice University Department of Electrical and Computer Engineering. (n.d.). \textit{ELEC 241: Fundamentals of electrical engineering I}. \url{https://www.ece.rice.edu/~dhj/courses/elec241/}

Stanford Engineering Everywhere. (n.d.). \textit{EE261: The Fourier transform and its applications}. Stanford University. \url{https://see.stanford.edu/Course/EE261}

\paragraph{14.3.4 หมายเหตุด้านการอ้างอิง}

\begin{itemize}
\item
  รายการอ้างอิงจัดรูปแบบตามแนว APA 7 โดยปรับให้เหมาะกับเอกสารการสอนภาษาไทย
\item
  วันที่เข้าถึงแหล่งออนไลน์สำหรับการจัดทำต้นฉบับนี้คือ 15 กรกฎาคม 2569
\item
  ก่อนเผยแพร่ฉบับพิมพ์ ควรตรวจสอบนโยบายลิงก์ การเปลี่ยนแปลงหน้าเว็บ และรายละเอียดฉบับพิมพ์ของหนังสืออีกครั้ง
\end{itemize}

%\par\noindent\rule{\textwidth}{0.4pt}\par

\subsection{ภาคผนวก ก สรุปสูตรสำคัญ}

\subsubsection{ก.1 อนุกรมเลขคณิต}

\[
a_n=a_1+(n-1)d
\]

\[
S_N=\frac{N}{2}[2a_1+(N-1)d]=\frac{N}{2}(a_1+a_N)
\]

\subsubsection{ก.2 อนุกรมเรขาคณิต}

\[
a_n=a_1r^{n-1}
\]

\[
S_N=a_1\frac{1-r^N}{1-r},\qquad r\ne1
\]

\[
S_\infty=\frac{a_1}{1-r},\qquad |r|<1
\]

\subsubsection{ก.3 ความถี่ของฟังก์ชันเป็นคาบ}

\[
f_0=\frac{1}{T},\qquad \omega_0=\frac{2\pi}{T},\qquad f_n=nf_0
\]

\subsubsection{ก.4 อนุกรมฟูเรียร์คาบ \(T\)}

\[
f(t)=\frac{a_0}{2}+\sum_{n=1}^{\infty}[a_n\cos(n\omega_0t)+b_n\sin(n\omega_0t)]
\]

\[
a_0=\frac{2}{T}\int_{t_0}^{t_0+T}f(t)dt
\]

\[
a_n=\frac{2}{T}\int_{t_0}^{t_0+T}f(t)\cos(n\omega_0t)dt
\]

\[
b_n=\frac{2}{T}\int_{t_0}^{t_0+T}f(t)\sin(n\omega_0t)dt
\]

\subsubsection{ก.5 สมมาตร}

\begin{itemize}
\item
  ฟังก์ชันคู่: \(b_n=0\)
\item
  ฟังก์ชันคี่: \(a_0=a_n=0\)
\item
  สมมาตรครึ่งคลื่น: DC และฮาร์มอนิกคู่เป็นศูนย์
\end{itemize}

\subsubsection{ก.6 ขนาดและเฟส}

\[
A_n=\sqrt{a_n^2+b_n^2}
\]

\[
\theta_n=\operatorname{atan2}(b_n,a_n)
\]

\[
a_n\cos(n\omega_0t)+b_n\sin(n\omega_0t)
=A_n\cos(n\omega_0t-\theta_n)
\]

\subsubsection{ก.7 รูปเชิงซ้อน}

\[
f(t)=\sum_{n=-\infty}^{\infty}C_ne^{jn\omega_0t}
\]

\[
C_n=\frac{1}{T}\int_{t_0}^{t_0+T}f(t)e^{-jn\omega_0t}dt
\]

\subsubsection{ก.8 พาร์เซวาล}

\[
\frac{1}{T}\int_{t_0}^{t_0+T}|f(t)|^2dt
=\frac{a_0^2}{4}+\frac{1}{2}\sum_{n=1}^{\infty}(a_n^2+b_n^2)
\]

%\par\noindent\rule{\textwidth}{0.4pt}\par

\subsection{ภาคผนวก ข แบบตรวจสอบก่อนส่งงานของนักศึกษา}

\begin{itemize}
\item
  {[} {]} ระบุคาบ \(T\) และ \(\omega_0\) ถูกต้อง
\item
  {[} {]} เลือกช่วงอินทิเกรตครบหนึ่งคาบ
\item
  {[} {]} เขียนฟังก์ชันเป็นช่วงครบถ้วน
\item
  {[} {]} ตรวจสมมาตรก่อนคำนวณ
\item
  {[} {]} แยกค่าเฉลี่ย \(a_0/2\) จาก \(a_0\)
\item
  {[} {]} ใช้สูตรสัมประสิทธิ์สอดคล้องกับนิยามอนุกรม
\item
  {[} {]} แสดงขั้นตอนอินทิเกรตและจัดรูปตาม \(n\)
\item
  {[} {]} ตรวจกรณี \(n\) คู่และคี่
\item
  {[} {]} ระบุหน่วยของสัมประสิทธิ์และความถี่
\item
  {[} {]} เขียนพจน์แรกเพื่อให้เห็นรูปแบบ
\item
  {[} {]} ตรวจสอบคำตอบด้วยสมมาตร ค่าเฉลี่ย หรือค่าตัวอย่าง
\item
  {[} {]} ตีความ DC มูลฐาน ฮาร์มอนิก และสเปกตรัม
\end{itemize}
